% bismi ALLAH ARRAHMAN ARRAHIM
































































































































































\documentclass[11pt,a4paper]{article}
\usepackage{booktabs}

% -------------------------------------------------
% Encodage, langue, typographie
% -------------------------------------------------
\usepackage[french]{babel}
\usepackage{fontspec}
\usepackage{microtype}
\usepackage[a4paper,margin=2.5cm]{geometry}
\usepackage{hyphenat}


% -------------------------------------------------
% Mathématiques
% -------------------------------------------------
\usepackage{amsmath,amssymb,amsthm,mathtools,bm}
\usepackage{enumitem}

% -------------------------------------------------
% Liens
% -------------------------------------------------
\usepackage[hidelinks]{hyperref}

% -------------------------------------------------
% Police
% -------------------------------------------------
\setmainfont{TeX Gyre Pagella}

% -------------------------------------------------
% Environnements
% -------------------------------------------------
\newtheorem{definition}{Définition}[section]
\newtheorem{assumption}{Hypothèse}[section]
\newtheorem{proposition}{Proposition}[section]
\newtheorem{theorem}{Théorème}[section]
\newtheorem{corollary}{Corollaire}[section]
\newtheorem{remark}{Remarque}[section]
\newtheorem{lemma}{Lemme}[section]

% -------------------------------------------------
% Opérateurs
% -------------------------------------------------
\DeclareMathOperator{\prox}{prox}
\DeclareMathOperator{\argmin}{argmin}
\DeclareMathOperator{\argmax}{argmax}
\DeclareMathOperator{\supp}{supp}
\DeclareMathOperator{\dist}{dist}
\DeclareMathOperator{\KL}{KL}
\DeclareMathOperator{\Id}{Id}
\DeclareMathOperator{\diag}{diag}
\DeclareMathOperator{\softmax}{softmax}
\DeclareMathOperator{\conv}{conv}
\DeclareMathOperator{\tr}{tr}

% -------------------------------------------------
% Raccourcis
% -------------------------------------------------
\newcommand{\X}{\mathcal X}
\newcommand{\C}{\mathcal C}
\newcommand{\D}{\mathcal D}
\newcommand{\E}{\mathcal E}
\newcommand{\Z}{\mathcal Z}
\newcommand{\Hh}{\mathcal H}
\newcommand{\I}{\mathcal I}
\newcommand{\T}{\mathcal T}
\newcommand{\W}{\mathcal W}
\newcommand{\N}{\mathcal N}
\newcommand{\R}{\mathbb R}
\newcommand{\Pcal}{\mathcal P}
\newcommand{\Qcal}{\mathcal Q}
\newcommand{\Gcal}{\mathcal G}
\newcommand{\Mcal}{\mathcal M}
\newcommand{\1}{\mathbf 1}

\title{PCH-ESM\\
\large Purification Contextuelle Hiérarchique et Espaces Sémantiques Modulaires}
\author{}
\date{}

\begin{document}
\emergencystretch=3em
\maketitle

\begin{abstract}
Ce document formalise une architecture hiérarchique pour modèles séquentiels et systèmes de langage, fondée sur quatre briques explicites et compatibles :
(i) une purification récursive d'une représentation dans un contexte donné ;
(ii) une décomposition orthogonale et déterministe d'un espace latent global en sous-espaces spécialisés obtenus par construction spectrale ;
(iii) une spécialisation locale certifiée par une fonctionnelle explicite d'apprentissage, puis un opérateur de coordination pondéré à retour ;
(iv) une dynamique Multi-PH pour un état latent hiérarchique interprété comme état du monde, avec croyance et incertitude explicites.
Chaque objet central est défini par projection, distance, argmin, ou opération spectrale. Aucune notion structurante n'est laissée implicite.
\end{abstract}

\tableofcontents

% =================================================
\section{Cadre général}
% =================================================


PCH-ESM formalise une architecture hiérarchique de représentation et de coordination fondée sur quatre idées centrales.

La première est la purification contextuelle hiérarchique : une représentation latente peut être ramenée vers un état stable par une minimisation variationnelle explicite, d’abord dans le cadre PH, puis dans le cadre contextuel PCH, et enfin dans une version dynamique Multi-PH munie d’un état de croyance et d’une incertitude.

La deuxième est la spécialisation orthogonale de l’espace latent : au lieu d’un espace monolithique, la théorie construit une décomposition spectrale en sous-espaces spécialisés, avec une partition canonique optimale lorsque le spectre est simple, et une formulation déterministe en termes de projecteurs spectraux.

La troisième est la coordination globale par collage pondéré : les sorties locales certifiées sont réconciliées par un opérateur de fusion explicite, optimal sous bruit hétéroscédastique au sens inverse-variance, et compatible avec une reconstruction globale contrôlée.

La quatrième est l’unification de ces briques dans une théorie complète de l’architecture : spécialisation, collage, découplage inter-modules, dynamique du monde et compromis de capacité sont formulés par des objectifs explicites, des minimisations convexes, des projections orthogonales et des critères d’optimalité démontrables.

L’objectif de ce document est de donner une formalisation unifiée, déterministe et analytiquement vérifiable de cette architecture, en séparant clairement les objets continus, les décisions spectrales, les mécanismes de coordination et la dynamique d’état.

\subsection*{Séparation des objets continus et discrets}

La théorie sépare strictement deux niveaux :

\begin{itemize}[leftmargin=2em]
    \item le niveau \emph{continu}, où l'on purifie des états latents par minimisation d'énergies ;
    \item le niveau \emph{discret}, où l'on choisit des espaces actifs par structure spectrale et par règles de départage fixées.
\end{itemize}

Cette séparation évite de mélanger dans un seul objet des opérations d'optimisation convexe et des décisions discrètes.

% =================================================
\section{Hypothèses et notations globales}
% =================================================

\begin{assumption}[Espace latent fini-dimensionnel]
On fixe un espace euclidien réel \(\Hh=\R^d\).  
Les représentations latentes vivent dans \(\Hh\) ou dans \(\Hh^T=\R^{dT}\) lorsqu'il s'agit de fenêtres séquentielles de longueur \(T\).
\end{assumption}

\begin{assumption}[Fenêtres et normes]
Pour une fenêtre \(u=(u_1,\dots,u_T)\in \Hh^T\), on considère la norme produit
\[
\|u\|^2 = \sum_{t=1}^T \|u_t\|^2.
\]
\end{assumption}

\begin{assumption}[Base de calibration]
On fixe un ensemble de calibration fini
\[
\D_0=\{x_1,\dots,x_N\}\subset \X
\]
muni de labels contextuels \(c_j\in\C\).  
On fixe aussi un encodeur de calibration \(E_{\theta_0}:\X\to \Hh^T\), tenu fixe lors de la construction des ensembles contextuels.
\end{assumption}

\begin{assumption}[Fidélité et régularité]
Les coefficients \(\lambda_{\mathrm{ch}},\lambda_{\mathrm{reg}},\lambda_{\mathrm{ctx}},\lambda_{i,\mathrm{reg}},\lambda_{i,\mathrm{ctx}},\lambda_{\mathrm{fb}} \\\lambda_{\mathrm{dyn}},\lambda_{\mathrm{sm}},\lambda_{\mathrm{spec}},\lambda_{\mathrm{glue}},\lambda_{\mathrm{coord}},\lambda_{\mathrm{task}}\) sont strictement positifs.
\end{assumption}

% =================================================
\section{PH : purification hiérarchique}
% =================================================

Dans cette théorie, PH est défini comme un opérateur de purification associé à une pénalisation explicite de chaîne et de norme.

\subsection{Énergie de PH}

\begin{definition}[Partie régulière de PH]
On définit
\[
\Phi_{\mathrm{PH}}(u)
=
\lambda_{\mathrm{ch}}\,\Omega_{\mathrm{ch}}(u)
+
\lambda_{\mathrm{reg}}\|u\|^2,
\qquad
\Omega_{\mathrm{ch}}(u)=\sum_{t=2}^T\|u_t-u_{t-1}\|^2.
\]
\end{definition}

\begin{definition}[Opérateur PH]
L'opérateur de purification hiérarchique est défini par
\[
\Pcal^{\mathrm{PH}}(w)
=
\argmin_{u\in\Hh^T}
\left\{
\frac12\|u-w\|^2+\Phi_{\mathrm{PH}}(u)
\right\}.
\]
\end{definition}

\begin{proposition}[Existence et unicité de PH]
Si \(\lambda_{\mathrm{reg}}>0\), alors \(\Phi_{\mathrm{PH}}\) est fortement convexe.  
Par conséquent, pour tout \(w\in\Hh^T\), l'opérateur \(\Pcal^{\mathrm{PH}}(w)\) existe et est unique.
\end{proposition}

\begin{proof}
Le terme \(u\mapsto \lambda_{\mathrm{reg}}\|u\|^2\) est fortement convexe.  
La somme avec la fonction convexe \(\lambda_{\mathrm{ch}}\Omega_{\mathrm{ch}}(u)\) demeure fortement convexe.  
L'ajout du terme \(\frac12\|u-w\|^2\) conserve cette propriété.  
Une fonction fortement convexe admet un minimiseur unique.
\end{proof}

%%%%%%%%%%%%%%%%%%%%%%%%%%%%%%%%%%%%%%%%%%%%%%%%%%%%%%%%%%%%%%%%%%%
%\begin{remark}
%Le schéma bidirectionnel du document initial peut être interprété comme un solveur numérique %du problème ci-dessus.  
%Il ne définit pas un objet concurrent : il réalise approximativement l'opérateur PH.
%\end{remark}
%%%%%%%%%%%%%%%%%%%%%%%%%%%%%%%%%%%%%%%%%%%%%%%%%%%%%%%%%%%%%%%%%%%

\subsection{Stabilité du mapping de PH}

\begin{theorem}[Lipschitz de l'opérateur PH]
Pour \(\lambda_{\mathrm{reg}}>0\), l'application
\[
w\longmapsto \Pcal^{\mathrm{PH}}(w)
\]
est bien définie et lipschitzienne. Plus précisément, il existe une constante \(L_{\mathrm{PH}}\in(0,1]\) telle que
\[
\|\Pcal^{\mathrm{PH}}(w_1)-\Pcal^{\mathrm{PH}}(w_2)\|
\le
L_{\mathrm{PH}}\|w_1-w_2\|
\qquad
\text{pour tous } w_1,w_2\in\Hh^T.
\]
\end{theorem}

\begin{proof}
On pose \(\Phi_{\mathrm{PH}}=\lambda_{\mathrm{ch}}\Omega_{\mathrm{ch}}+\lambda_{\mathrm{reg}}\|\cdot\|^2\).  
Comme \(\Phi_{\mathrm{PH}}\) est fortement convexe, son opérateur proximal \(\prox_{\Phi_{\mathrm{PH}}}\) est contractant.  
Or \(\Pcal^{\mathrm{PH}}=\prox_{\Phi_{\mathrm{PH}}}\).  
Le résultat suit.
\end{proof}

\subsection{Trajectoire proximale de PH}

\begin{definition}[Itération proximale de PH]
Pour \(\eta>0\), on définit
\[
\mathcal T^{\mathrm{PH}}_{\eta}(z;w)
=
\argmin_{u\in\Hh^T}
\left\{
\frac12\|u-w\|^2
+\Phi_{\mathrm{PH}}(u)
+\frac{1}{2\eta}\|u-z\|^2
\right\}.
\]
La trajectoire proximale est
\[
z^{(0)}=w,
\qquad
z^{(r+1)}=\mathcal T^{\mathrm{PH}}_{\eta}(z^{(r)};w),
\quad r\ge 0.
\]
\end{definition}

\begin{theorem}[Contraction de l'itération proximale PH]
Soit \(J_{\mathrm{PH}}(u;w)=\frac12\|u-w\|^2+\Phi_{\mathrm{PH}}(u)\).  
Alors \(J_{\mathrm{PH}}(\cdot;w)\) est fortement convexe, et pour tout \(\eta>0\), l'opérateur \(\mathcal T^{\mathrm{PH}}_{\eta}(\cdot;w)\) est une contraction.  
Par conséquent, la suite \(z^{(r)}\) converge linéairement vers l'unique point fixe \(\bar z_{\mathrm{PH}}(w)\), qui coïncide avec \(\Pcal^{\mathrm{PH}}(w)\).
\end{theorem}

\begin{proof}
Soient \(x=\mathcal T^{\mathrm{PH}}_{\eta}(a;w)\) et \(y=\mathcal T^{\mathrm{PH}}_{\eta}(b;w)\).  
Par optimalité,
\[
\frac{a-x}{\eta}\in \partial J_{\mathrm{PH}}(x;w),
\qquad
\frac{b-y}{\eta}\in \partial J_{\mathrm{PH}}(y;w).
\]
Comme \(J_{\mathrm{PH}}(\cdot;w)\) est fortement convexe, son sous-différentiel est fortement monotone.  
On obtient donc
\[
\langle x-y,a-b\rangle\ge (1+\eta\mu_{\mathrm{PH}})\|x-y\|^2
\]
pour un certain \(\mu_{\mathrm{PH}}>0\).  
Par Cauchy-Schwarz,
\[
\|x-y\|\le \frac{1}{1+\eta\mu_{\mathrm{PH}}}\|a-b\|.
\]
L'opérateur est contractant, donc la convergence suit par Banach.  
Le point fixe vérifie la condition d'optimalité du minimiseur.
\end{proof}

\begin{proposition}[Décroissance de l'énergie en PH]
La trajectoire proximale vérifie
\[
J_{\mathrm{PH}}(z^{(r+1)};w)
+
\frac{1}{2\eta}\|z^{(r+1)}-z^{(r)}\|^2
\le
J_{\mathrm{PH}}(z^{(r)};w),
\]
d'où la décroissance monotone de \(J_{\mathrm{PH}}(z^{(r)};w)\).
\end{proposition}

\begin{proof}
Par définition de \(z^{(r+1)}\) comme minimiseur du problème proximal, en comparant à \(u=z^{(r)}\), on obtient immédiatement l'inégalité.
\end{proof}

% =================================================
\section{De PH vers PCH : purification contextuelle hiérarchique}
% =================================================

Cette section généralise PH en introduisant explicitement un contexte admissible.

\subsection{Construction explicite des ensembles contextuels}

\begin{definition}[Échantillon contextuel]
Pour chaque contexte \(c\in\C\), on note
\[
\D_c=\{x_j\in\D_0 : c_j=c\}.
\]
\end{definition}

\begin{definition}[Ensemble admissible contextuel]
Pour chaque contexte \(c\), on définit
\[
\Mcal_c
=
\conv\{E_{\theta_0}(x): x\in\D_c\}\subseteq \Hh^T.
\]
Comme \(\D_c\) est fini, \(\Mcal_c\) est compact et convexe.
\end{definition}

\begin{lemma}[Projection sur l'ensemble contextuel]
Pour tout \(u\in\Hh^T\) et tout \(c\in\C\), la projection orthogonale
\[
\Pi_{\Mcal_c}(u)=\argmin_{v\in\Mcal_c}\|u-v\|
\]
existe et est unique.
\end{lemma}

\begin{proof}
En dimension finie, \(\Mcal_c\) est non vide, fermé et convexe.  
La projection orthogonale sur un tel ensemble existe et est unique.
\end{proof}
%%%%%%%%%%%%%%%%%%%%%%%%%%%%%%%%%%%%%%%%%%%%%%%%%%%%%%%%%%%%%%%%%
%\begin{remark}
%Cette définition ferme la difficulté pratique restante : le contexte n'est pas un prototype %caché.  
%Il est une région admissible dérivée d'un corpus calibré et d'un encodeur figé.
%\end{remark}
%%%%%%%%%%%%%%%%%%%%%%%%%%%%%%%%%%%%%%%%%%%%%%%%%%%%%%%%%%%%%%%%%

\subsection{Fonctionnelle contextuelle}

\begin{definition}[Partie contextuelle régulière]
On définit
\[
\Phi_{\mathrm{PCH},c}(u)
=
\lambda_{\mathrm{ctx}}\,\dist(u,\Mcal_c)^2
+
\lambda_{\mathrm{ch}}\,\Omega_{\mathrm{ch}}(u)
+
\lambda_{\mathrm{reg}}\|u\|^2.
\]
\end{definition}

\begin{definition}[Opérateur PCH]
L'opérateur de purification contextuelle hiérarchique est défini par
\[
\Pcal^{\mathrm{PCH}}(w;c)
=
\argmin_{u\in\Hh^T}
\left\{
\frac12\|u-w\|^2+\Phi_{\mathrm{PCH},c}(u)
\right\}.
\]
\end{definition}

\begin{proposition}[Existence et unicité de PCH]
Si \(\lambda_{\mathrm{reg}}>0\), alors \(\Phi_{\mathrm{PCH},c}\) est fortement convexe pour tout \(c\).  
Par conséquent, pour tout \(w\) et tout \(c\), \(\Pcal^{\mathrm{PCH}}(w;c)\) existe et est unique.
\end{proposition}

\begin{proof}
Le terme \(u\mapsto \lambda_{\mathrm{reg}}\|u\|^2\) est fortement convexe.  
Les termes \(\lambda_{\mathrm{ctx}}\dist(u,\Mcal_c)^2\) et \(\lambda_{\mathrm{ch}}\Omega_{\mathrm{ch}}(u)\) sont convexes.  
La somme reste fortement convexe, donc admet un minimiseur unique.
\end{proof}

\subsection{Stabilité du mapping contextuel}

\begin{theorem}[Lipschitz de l'opérateur PCH]
Pour chaque contexte \(c\), l'application
\[
w\longmapsto \Pcal^{\mathrm{PCH}}(w;c)
\]
est bien définie et lipschitzienne.  
En particulier, il existe une constante \(L_{\mathrm{PCH}}\in(0,1]\) telle que
\[
\|\Pcal^{\mathrm{PCH}}(w_1;c)-\Pcal^{\mathrm{PCH}}(w_2;c)\|
\le
L_{\mathrm{PCH}}\|w_1-w_2\|
\]
pour tous \(w_1,w_2\in\Hh^T\).
\end{theorem}

\begin{proof}
On écrit
\[
\Pcal^{\mathrm{PCH}}(w;c)=\prox_{\Phi_{\mathrm{PCH},c}}(w).
\]
Comme \(\Phi_{\mathrm{PCH},c}\) est fortement convexe, son proximal est contractant.  
La conclusion suit.
\end{proof}

\subsection{Itération proximale de PCH}

\begin{definition}[Itération proximale de PCH]
Pour \(\eta>0\), on définit
\[
\mathcal T^{\mathrm{PCH}}_{\eta}(z;w,c)
=
\argmin_{u\in\Hh^T}
\left\{
\frac12\|u-w\|^2+\Phi_{\mathrm{PCH},c}(u)+\frac{1}{2\eta}\|u-z\|^2
\right\}.
\]
La trajectoire proximale est
\[
u^{(0)}=w,
\qquad
u^{(r+1)}=\mathcal T^{\mathrm{PCH}}_{\eta}(u^{(r)};w,c),
\quad r\ge 0.
\]
\end{definition}

\begin{theorem}[Contraction de l'itération proximale PCH]
Soit \(J_{\mathrm{PCH}}(u;w,c)=\frac12\|u-w\|^2+\Phi_{\mathrm{PCH},c}(u)\).  
Alors \(J_{\mathrm{PCH}}(\cdot;w,c)\) est fortement convexe, et pour tout \(\eta>0\), l'opérateur \(\mathcal T^{\mathrm{PCH}}_{\eta}(\cdot;w,c)\) est une contraction.  
Par conséquent, la suite \(u^{(r)}\) converge linéairement vers l'unique point fixe \(\bar z_{\mathrm{PCH}}(w,c)\), qui coïncide avec \(\Pcal^{\mathrm{PCH}}(w;c)\).
\end{theorem}

\begin{proof}
La preuve est identique à celle de PH, en remplaçant \(J_{\mathrm{PH}}\) par \(J_{\mathrm{PCH}}\).  
La forte convexité donne la forte monotonie du sous-différentiel, puis la contraction.  
Banach donne la convergence vers un unique point fixe.  
Ce point fixe vérifie la condition d'optimalité du minimiseur.
\end{proof}

\begin{proposition}[Décroissance de l'énergie en PCH]
La trajectoire proximale vérifie
\[
J_{\mathrm{PCH}}(u^{(r+1)};w,c)
+
\frac{1}{2\eta}\|u^{(r+1)}-u^{(r)}\|^2
\le
J_{\mathrm{PCH}}(u^{(r)};w,c),
\]
d'où la décroissance monotone de \(J_{\mathrm{PCH}}(u^{(r)};w,c)\).
\end{proposition}

\begin{proof}
Même argument que pour PH : on compare le minimum proximal à la valeur en \(u^{(r)}\).
\end{proof}

\begin{theorem}[PH est un cas limite de PCH]
Supposons que le contexte soit trivial au sens suivant :
\[
\Mcal_c=\Hh^T
\quad\text{pour tout } c\in\C.
\]
Alors
\[
\dist(u,\Mcal_c)=0
\quad\text{pour tout } u\in\Hh^T,
\]
et
\[
\Pcal^{\mathrm{PCH}}(w;c)=\Pcal^{\mathrm{PH}}(w).
\]
Par conséquent, les trajectoires proximales coïncident.
\end{theorem}

\begin{proof}
Si \(\Mcal_c=\Hh^T\), le terme contextuel est identiquement nul.  
La fonctionnelle PCH se réduit alors exactement à PH, et les minimisateurs coïncident.
\end{proof}

\subsection{Stabilité contextuelle canonique}

\begin{definition}[Instabilité contextuelle fondamentale]
On définit
\[
\mathcal I(z;w,c)
=
J_{\mathrm{PCH}}(z;w,c)-\inf_{u\in\Hh^T}J_{\mathrm{PCH}}(u;w,c).
\]
\end{definition}

\begin{definition}[Stabilité contextuelle fondamentale]
On définit le score canonique de stabilité par
\[
S_C^\star(z;w,c)=\frac{1}{1+\mathcal I(z;w,c)}.
\]
\end{definition}

\begin{proposition}[Bornes]
Pour tout \(z,w,c\),
\[
0<S_C^\star(z;w,c)\le 1.
\]
De plus, \(S_C^\star=1\) si et seulement si \(z\) minimise \(J_{\mathrm{PCH}}(\cdot;w,c)\).
\end{proposition}

\begin{proof}
Par définition, \(\mathcal I\ge 0\).  
Le reste est immédiat.
\end{proof}

\begin{definition}[Gain de stabilité]
On définit
\[
\mathsf{Gain}(w,c)
=
J_{\mathrm{PCH}}\!\bigl(w;w,c\bigr)
-
J_{\mathrm{PCH}}\!\bigl(\Pcal^{\mathrm{PCH}}(w;c);w,c\bigr).
\]
\end{definition}

\begin{proposition}[Gain de stabilité non négatif]
On a \(\mathsf{Gain}(w,c)\ge 0\) pour tout \(w,c\).
\end{proposition}

\begin{proof}
Par définition, \(\Pcal^{\mathrm{PCH}}(w;c)\) minimise \(J_{\mathrm{PCH}}(\cdot;w,c)\).  
Donc
\[
J_{\mathrm{PCH}}(\Pcal^{\mathrm{PCH}}(w;c);w,c)\le J_{\mathrm{PCH}}(w;w,c),
\]
d'où \(\mathsf{Gain}(w,c)\ge 0\).
\end{proof}

\begin{remark}
Le score \(S_C^\star\) est une stabilité relative à la fonctionnelle PCH elle-même.  
Il ne prétend pas mesurer une vérité universelle, mais une stabilité certifiée dans le cadre formel défini ici.
\end{remark}

% =================================================
\section{PCH-ESM : base spectrale, spécialisation certifiée et apprentissage}
% =================================================

Cette section formalise la décentralisation de l'espace latent, la spécialisation locale, et la construction déterministe des sous-espaces.  
Le point essentiel est que la spécialisation n'est plus seulement une intuition : elle est définie par un objectif explicite, des certificats de marge, et une condition d'identifiabilité.

\subsection{Construction spectrale déterministe}

\begin{definition}[Covariance de calibration]
Soient \(z_j=E_{\theta_0}(x_j)\in\Hh^T\) les latents de calibration.  
On définit la moyenne
\[
\bar z=\frac1N\sum_{j=1}^N z_j
\]
et la covariance empirique symétrique
\[
\Sigma=\frac1N\sum_{j=1}^N (z_j-\bar z)(z_j-\bar z)^\top\in\R^{dT\times dT}.
\]
\end{definition}

\begin{definition}[Projecteurs spectraux]
Par le théorème spectral, il existe une famille finie de projecteurs orthogonaux \((P_1,\dots,P_m)\) associés aux sous-espaces propres de \(\Sigma\).  
Lorsque certaines valeurs propres sont multiples, on travaille avec les projecteurs spectraux et non avec une base arbitraire ; la construction reste donc déterministe.
\end{definition}

\begin{definition}[Partition spectrale]
On fixe une partition déterministe de l'ensemble des indices spectraux \(\{1,\dots,m\}\) en blocs disjoints \(B_0,\dots,B_n\).  
On définit alors
\[
\Hh_i=\bigoplus_{k\in B_i}\mathrm{Im}(P_k),\qquad i=0,\dots,n.
\]
\end{definition}

\begin{definition}[Projecteurs sur les sous-espaces spécialisés]
On note \(\Pi_i:\Hh^T\to\Hh_i\) la projection orthogonale sur \(\Hh_i\).  
Elle est linéaire, idempotente et auto-adjointe.
\end{definition}

\begin{proposition}[Décomposition orthogonale unique]
On a
\[
\Hh^T=\Hh_0\oplus\Hh_1\oplus\cdots\oplus\Hh_n.
\]
Pour tout \(z\in\Hh^T\), il existe une unique écriture
\[
z=z_0+\cdots+z_n,
\qquad
z_i\in\Hh_i.
\]
De plus,
\[
\|z\|^2=\|z_0\|^2+\cdots+\|z_n\|^2.
\]
\end{proposition}

\begin{proof}
Les sous-espaces \(\Hh_i\) sont construits à partir de projecteurs spectraux orthogonaux.  
Ils sont donc deux à deux orthogonaux et leur somme directe est l'espace total.  
L'unicité et la formule de Pythagore suivent immédiatement.
\end{proof}

\begin{remark}
La base de spécialisation n'est pas postulée comme une boîte noire : elle est obtenue par projecteurs spectraux.  
La détermination est explicite et ne dépend pas d'un choix arbitraire de base dans les sous-espaces propres dégénérés.
\end{remark}

\subsection{Signature de spécialité certifiée}

\begin{definition}[Signature probabiliste complète]
Pour \(z\in\Hh^T\), on définit la signature
\[
\sigma(z)=\bigl(\sigma_0(z),\sigma_1(z),\dots,\sigma_n(z)\bigr)
\]
par
\[
\sigma_i(z)
=
\frac{\|\Pi_i z\|^2+\varepsilon}{\|z\|^2+(n+1)\varepsilon},
\qquad i=0,\dots,n,
\]
où \(\varepsilon>0\) évite la division par zéro.
\end{definition}

\begin{proposition}[Normalisation]
Pour tout \(z\in\Hh^T\),
\[
\sum_{i=0}^n \sigma_i(z)=1.
\]
\end{proposition}

\begin{proof}
Par définition,
\[
\sum_{i=0}^n(\|\Pi_i z\|^2+\varepsilon)
=
\|z\|^2+(n+1)\varepsilon.
\]
La somme des composantes est donc exactement \(1\).
\end{proof}

\begin{remark}
La signature n'est pas une étiquette postulée.  
Elle est dérivée de la géométrie effective du latent et peut être utilisée comme variable de spécialisation certifiée.
\end{remark}

\subsection{Objectif de spécialisation}

\begin{assumption}[Séparabilité des domaines]
Pour chaque domaine \(D_i\), il existe une marge \(\delta_i>0\) telle que, sur les données de calibration du domaine \(D_i\),
\[
\mathbb E_{x\sim D_i}\bigl[\sigma_i(E_{\theta_0}(x))\bigr]
\ge
\max_{j\neq i}\mathbb E_{x\sim D_i}\bigl[\sigma_j(E_{\theta_0}(x))\bigr]
+\delta_i.
\]
\end{assumption}

\begin{definition}[Fonctionnelle de spécialisation]
Pour des projecteurs \(\mathbf \Pi=(\Pi_0,\dots,\Pi_n)\) et des têtes locales \(\mathbf h=(h_1,\dots,h_n)\), on définit
\[
\mathcal J_{\mathrm{spec}}(\mathbf\Pi,\mathbf h)
=
\sum_{i=1}^n
\mathbb E_{(x,y)\sim D_i}
\Bigl[
\ell_i\bigl(h_i(\Pi_i E_{\theta_0}(x)),y\bigr)
\Bigr]
+
\lambda_{\mathrm{leak}}
\sum_{i\neq j}
\mathbb E_{x\sim D_i}\|\Pi_j E_{\theta_0}(x)\|^2
+
\lambda_{\mathrm{bal}}
\sum_{i=0}^n
\bigl(\operatorname{tr}(\Pi_i)-d_i\bigr)^2.
\]
Ici \(\ell_i\) est une perte locale convexe pour la tâche \(i\), et les têtes \(h_i\) sont prises dans une classe bornée de maps linéaires ou affines.
\end{definition}

\begin{proposition}[Existence d'une spécialisation optimale]
L'ensemble des tuples de projecteurs orthogonaux \((\Pi_0,\dots,\Pi_n)\) de rangs fixés est compact, et la classe bornée des têtes \(\mathbf h\) est compacte en dimension finie.  
Par conséquent, \(\mathcal J_{\mathrm{spec}}\) admet au moins un minimiseur global.
\end{proposition}

\begin{proof}
En dimension finie, l'ensemble des projecteurs orthogonaux de rang fixé s'identifie à un produit de variétés de Grassmann compactes.  
La classe bornée des têtes est compacte par Heine-Borel.  
La fonction \(\mathcal J_{\mathrm{spec}}\) est continue, donc elle atteint son minimum.
\end{proof}

\begin{theorem}[Identifiabilité de la spécialisation]
Sous l'hypothèse de séparabilité des domaines, si un tuple \((\mathbf \Pi^\star,\mathbf h^\star)\) minimise \(\mathcal J_{\mathrm{spec}}\), alors chaque domaine \(D_i\) est associé de manière identifiable à son bloc \(\Hh_i\) au sens suivant :
pour tout \(j\neq i\),
\[
\mathbb E_{x\sim D_i}\bigl[\sigma_i(\Pi_i^\star E_{\theta_0}(x))\bigr]
>
\mathbb E_{x\sim D_i}\bigl[\sigma_j(\Pi_j^\star E_{\theta_0}(x))\bigr].
\]
En particulier, l'attribution des domaines aux blocs spécialisés est unique à permutation interne près dans les sous-espaces spectraux dégénérés.
\end{theorem}

\begin{proof}
La marge \(\delta_i>0\) impose une séparation stricte entre le score du bon bloc et ceux des blocs concurrents.  
Si un minimiseur violait cette séparation, il augmenterait nécessairement le terme de fuite \(\mathcal J_{\mathrm{spec}}\) d'au moins une quantité strictement positive.  
Ainsi, tout minimiseur respecte la séparation de marge et l'association domaine-bloc est identifiable.  
L'unicité n'est perdue qu'à l'intérieur des sous-espaces spectraux dégénérés, où la théorie travaille avec des projecteurs et non avec une base arbitraire.
\end{proof}

\begin{definition}[Apprentissage par alternance]
On considère l'itération bloc-coordinate déterministe
\[
\Pi^{(k+1)}
=
\argmin_{\Pi\in\mathfrak P}\mathcal J_{\mathrm{spec}}(\Pi,h^{(k)}),
\qquad
h^{(k+1)}
=
\argmin_{h\in\mathfrak H}\mathcal J_{\mathrm{spec}}(\Pi^{(k+1)},h),
\]
où \(\mathfrak P\) désigne l'ensemble admissible des partitions spectrales et \(\mathfrak H\) l'ensemble admissible des têtes bornées.
\end{definition}

\begin{proposition}[Décroissance de l'apprentissage]
La suite
\[
\mathcal J_{\mathrm{spec}}(\Pi^{(k)},h^{(k)})
\]
est décroissante.  
Toute suite d'itérations possède au moins un point d'accumulation, et tout point d'accumulation est stationnaire pour la fonctionnelle de spécialisation.
\end{proposition}

\begin{proof}
À chaque étape, on minimise exactement l'objectif sur l'un des blocs de variables.  
La valeur ne peut donc qu'être non croissante.  
La compacité des ensembles admissibles assure l'existence de points d'accumulation.  
Par continuité de l'objectif, tout point limite satisfait la condition stationnaire usuelle d'un minimum sur ensemble compact.
\end{proof}

\subsection{Contextes locaux spécialisés}

\begin{definition}[Contexte local]
Pour chaque contexte global \(c\in\C\), on considère des ensembles admissibles locaux
\[
\Mcal_{c,i}=\Pi_i(\Mcal_c)\subseteq \Hh_i,
\qquad i=0,\dots,n.
\]
Comme \(\Mcal_c\) est compact convexe et \(\Pi_i\) est linéaire, chaque \(\Mcal_{c,i}\) est compact convexe.
\end{definition}

\begin{assumption}[Admissibles locaux]
Pour chaque \(i\in\{0,\dots,n\}\), \(\Mcal_{c,i}\subseteq \Hh_i\) est non vide, fermé et convexe.
\end{assumption}

\begin{definition}[Opérateur local de purification]
Pour \(w_i\in\Hh_i\), on définit
\[
\Phi_i(w_i;c)
=
\argmin_{u_i\in\Hh_i}
\left\{
\frac12\|u_i-w_i\|^2
+
\lambda_{i,\mathrm{ctx}}\dist(u_i,\Mcal_{c,i})^2
+
\lambda_{i,\mathrm{reg}}\|u_i\|^2
\right\}.
\]
\end{definition}

\begin{proposition}[Unicité locale]
Si \(\lambda_{i,\mathrm{reg}}>0\), alors \(\Phi_i(\cdot;c)\) est bien définie et unique pour tout \(c\).
\end{proposition}

\begin{proof}
La fonction à minimiser est fortement convexe par la présence du terme quadratique \(\|u_i-w_i\|^2\) et du terme \(\lambda_{i,\mathrm{reg}}\|u_i\|^2\).  
Le minimiseur est donc unique.
\end{proof}

\begin{definition}[Sortie locale activée]
Pour une entrée locale \(w_i\in\Hh_i\), on pose
\[
u_i^\circ(w_i;c)=\Phi_i(w_i;c),
\]
puis le score de fiabilité et d'activation
\[
\rho_i(w,c)
=
\exp\!\bigl(-\|u_0^{(i)}(w,c)-\Pi_0 w\|^2\bigr)\,
\frac{\|\Pi_i w\|^2+\varepsilon}{\|w\|^2+(n+1)\varepsilon}\,
\frac{1}{1+\|u_i^\circ(\Pi_i w;c)-\Pi_i w\|^2}.
\]
Enfin, la sortie locale activée est
\[
\Psi_i(w,c)
=
\rho_i(w,c)\,u_i^\circ(\Pi_i w;c)
+
\bigl(1-\rho_i(w,c)\bigr)\Pi_i w.
\]
\end{definition}

%%%%%%%%%%%%%%%%%%%%%%%%%%%%%%%%%%%%%%%%%%%%%%%%%%%%%%%%%%%%%%%%%%%%%%%
%\begin{remark}
%Le score \(\rho_i\) n'est pas décoratif : il pondère directement la part corrigée par %l'opérateur local, puis il intervient dans le collage global.  
%Il mesure simultanément la proximité au sous-espace partagé, la concentration sur le bloc \  %(i\), et la stabilité de la sortie locale.
%\end{remark}
%%%%%%%%%%%%%%%%%%%%%%%%%%%%%%%%%%%%%%%%%%%%%%%%%%%%%%%%%%%%%%%%%%%%%%%

\begin{definition}[Certificat local]
Pour une entrée \(w\in\Hh^T\) et un contexte \(c\in\C\), le certificat local de l'espace \(i\) est le couple
\[
\mathfrak c_i(w,c)=\bigl(u_0^{(i)}(w,c),\Psi_i(w,c)\bigr)\in \Hh_0\oplus \Hh_i.
\]
\end{definition}

\begin{definition}[Sous-espace certifié local]
Pour chaque \(i\in\{1,\dots,n\}\), on définit
\[
\mathcal K_{c,i}
=
\{(\Pi_0 u,\Pi_i u): u\in \Mcal_c\}
\subseteq \Hh_0\oplus \Hh_i.
\]
Comme \(\Mcal_c\) est compact et convexe, \(\mathcal K_{c,i}\) est compact et convexe.
\end{definition}

\begin{proposition}[Existence et unicité du certificat local]
Si \(\lambda_{i,\mathrm{reg}}>0\), alors le certificat local construit à partir de la minimisation sur \(\mathcal K_{c,i}\) existe et est unique pour tout \(w\) et tout \(c\).
\end{proposition}

\begin{proof}
La fonction à minimiser est fortement convexe sur le sous-ensemble convexe compact \(\mathcal K_{c,i}\).  
Un minimiseur unique existe donc.
\end{proof}

\begin{definition}[Fusion globale]
On définit
\[
\Gamma(z_0,z_1,\dots,z_n)=z_0+\sum_{i=1}^n z_i.
\]
\end{definition}

\begin{proposition}[Cohérence de fusion]
Pour tout \(z\in\Hh^T\),
\[
\Gamma(\Pi_0 z,\Pi_1 z,\dots,\Pi_n z)=z.
\]
\end{proposition}

\begin{proof}
C'est la propriété de la décomposition orthogonale.
\end{proof}

% =================================================
\section{La fonction \(G\) : collage pondéré et coordination inter-espaces}
% =================================================

Dans cette théorie, \(G\) n'est pas une simple moyenne.  
C'est l'opérateur de collage pondéré à retour qui reconstruit une cohérence globale à partir de certificats locaux, et dont les poids dépendent explicitement de la fiabilité, de la spécialisation et de la stabilité.

\subsection{Fonctionnelle de coordination}

\begin{definition}[Fonctionnelle locale de coordination]
Pour une entrée \(w\in\Hh^T\) et un contexte \(c\in\C\), on définit
\[
\mathcal J_{\mathrm{G}}(v_0,v_1,\dots,v_n;w,c)
=
\sum_{i=1}^n \rho_i(w,c)\,\|v_0-u_0^{(i)}(w,c)\|^2
+
\lambda_{\mathrm{fb}}
\sum_{i=1}^n \rho_i(w,c)\,
\bigl\|v_i-\Psi_i(w,c)-\Pi_i\bigl(v_0-u_0^{(i)}(w,c)\bigr)\bigr\|^2.
\]
\end{definition}

\begin{definition}[Collage global \(G\)]
La fonction de coordination \(G\) est l'opérateur défini par
\[
G\bigl(\mathfrak c_1(w,c),\dots,\mathfrak c_n(w,c);w,c\bigr)
=
\bigl(\bar u_0(w,c),\bar u_1(w,c),\dots,\bar u_n(w,c)\bigr),
\]
où \((\bar u_0,\bar u_1,\dots,\bar u_n)\) est l'unique minimiseur de \(\mathcal J_{\mathrm{G}}(\cdot;w,c)\).
\end{definition}

\begin{proposition}[Forme explicite de la solution]
Le minimiseur de \(\mathcal J_{\mathrm{G}}\) vérifie
\[
\bar u_0(w,c)
=
\frac{\sum_{i=1}^n \rho_i(w,c)\,u_0^{(i)}(w,c)}{\sum_{i=1}^n \rho_i(w,c)},
\]
et
\[
\bar u_i(w,c)
=
\Psi_i(w,c)
+
\Pi_i\bigl(\bar u_0(w,c)-u_0^{(i)}(w,c)\bigr).
\]
\end{proposition}

\begin{proof}
Pour \(v_0\) fixé, chaque \(v_i\) minimise un terme quadratique strictement convexe, d'où
\[
v_i=\Psi_i+\Pi_i(v_0-u_0^{(i)}).
\]
En réinjectant dans l'objectif, on obtient une somme pondérée de carrés en \(v_0\), dont le minimiseur est la moyenne pondérée par les \(\rho_i\).  
La formule suit.
\end{proof}

\begin{remark}
La fonction \(G\) n'est donc ni une simple moyenne aveugle ni une boîte noire de distillation.  
Elle réalise un consensus pondéré sur la composante partagée et un retour correctif sur les composantes spécialisées.  
Les poids \(\rho_i\) contrôlent réellement l'influence de chaque espace.
\end{remark}

\begin{definition}[Disagreement de collage]
On définit le coût de désaccord partagé
\[
\mathcal E_{\mathrm{glue}}(w,c)
=
\sum_{i=1}^n
\rho_i(w,c)\,\|u_0^{(i)}(w,c)-\bar u_0(w,c)\|^2.
\]
\end{definition}

\begin{proposition}[Minimalité du consensus]
Le consensus \(\bar u_0(w,c)\) minimise \(\mathcal E_{\mathrm{glue}}(\cdot,c)\) parmi tous les choix possibles de composante partagée.
\end{proposition}

\begin{proof}
C'est la caractérisation standard du barycentre pondéré euclidien comme minimiseur d'une somme pondérée de carrés.
\end{proof}

\begin{proposition}[Robustesse du collage]
Si le certificat local \(k\) est perturbé en norme par \(\delta\), alors la perturbation induite sur le consensus partagé vérifie
\[
\|\Delta \bar u_0\|
\le
\frac{\rho_k(w,c)}{\sum_{i=1}^n \rho_i(w,c)}\,\delta.
\]
En particulier, si \(\rho_k(w,c)\) décroît, l'influence du certificat \(k\) décroît aussi.
\end{proposition}

\begin{proof}
La formule du consensus pondéré est linéaire en les certificats.  
La variation du consensus est donc bornée par le poids normalisé du certificat perturbé multiplié par la taille de la perturbation.
\end{proof}

\begin{theorem}[Exactitude du collage lorsque les certificats concordent]
Supposons que
\[
u_0^{(1)}(w,c)=\cdots=u_0^{(n)}(w,c)=u_0^\star(w,c).
\]
Alors
\[
\bar u_0(w,c)=u_0^\star(w,c),
\]
et le collage \(G\) restitue exactement la composante partagée certifiée.
\end{theorem}

\begin{proof}
Si tous les certificats partagés sont égaux, leur moyenne pondérée leur est égale.  
Le résultat suit immédiatement.
\end{proof}

\begin{definition}[Signature locale de coordination]
On définit la signature locale de coordination par
\[
\rho_i^{\mathrm{coord}}(w,c)
=
\exp\!\left(-\|u_0^{(i)}(w,c)-\bar u_0(w,c)\|^2\right).
\]
\end{definition}

\begin{theorem}[Reconnaissance de la faible fiabilité]
Si \(\rho_i(w,c)\to 0\), alors la contribution effective de l'espace \(i\) au collage global tend vers \(0\).  
En particulier, le certificat \(i\) devient asymptotiquement sans influence sur le collage global.
\end{theorem}

\begin{proof}
Par définition, les poids de \(G\) sont normalisés par la somme des fiabilités.  
Si \(\rho_i\) tend vers zéro alors que les autres fiabilités restent bornées strictement au-dessus de zéro, le poids normalisé de l'espace \(i\) tend vers zéro.
\end{proof}

% =================================================
\section{Multi-PH : état du monde, observations partielles et croyances}
% =================================================

On appelle Multi-PH la version où la purification hiérarchique ne sert plus seulement à construire une représentation utile pour une tâche, mais à modéliser l'évolution d'un état latent hiérarchique interprété comme état du monde.

\begin{remark}
Ici, le terme \emph{monde} ne signifie pas une causalité cachée supposée a priori.  
Il désigne l'état latent complet sur lequel s'exerce la dynamique variationnelle, avec observabilité partielle et incertitude explicitement modélisées.
\end{remark}

\subsection{Espace d'état}

\begin{assumption}[Décomposition hiérarchique de l'état]
On fixe un espace d'état euclidien
\[
\Hh_\Omega=\Hh^{(0)}\oplus\Hh^{(1)}\oplus\cdots\oplus\Hh^{(L)}.
\]
On identifie
\[
\Hh^{(0)}=\Hh_{\mathrm{obs}},
\qquad
\Hh^{(1)}=\Hh_{\mathrm{act}},
\]
et les niveaux \(\Hh^{(2)},\dots,\Hh^{(L)}\) comme niveaux de mémoire ou d'abstraction.
\end{assumption}

\begin{definition}[État hiérarchique]
Un état du monde à l'instant \(t\) est un vecteur
\[
\Omega_t=(\omega_t^{(0)},\omega_t^{(1)},\dots,\omega_t^{(L)})\in\Hh_\Omega.
\]
\end{definition}

\begin{definition}[Projecteurs d'observation et d'action]
On définit
\[
R:\Hh_\Omega\to\Hh_{\mathrm{obs}},
\qquad
A:\Hh_\Omega\to\Hh_{\mathrm{act}}
\]
par
\[
R(\Omega)=\omega^{(0)},
\qquad
A(\Omega)=\omega^{(1)}.
\]
\end{definition}

\subsection{Belief state et observabilité partielle}

\begin{definition}[Belief state]
On appelle croyance au temps \(t\) le triplet
\[
\mathfrak B_t=(\Omega_t,Q_t,r_t),
\]
où \(Q_t\succ 0\) est une métrique de confiance et \(r_t\ge 0\) un rayon d'incertitude.
\end{definition}

\begin{definition}[Ensemble de croyance]
À partir d'une observation partielle \(o_t\in\Hh_{\mathrm{obs}}\) et d'une action \(a_t\in\Hh_{\mathrm{act}}\), on définit l'ensemble admissible
\[
\mathcal B_t
=
\Bigl\{
\Omega\in\Hh_\Omega:
\|R\Omega-o_t\|_{W_t}\le \varepsilon_t,\;
\|A\Omega-a_t\|_{V_t}\le \eta_t,\;
\|\Omega-\Omega_t\|_{Q_t}\le r_t
\Bigr\},
\]
où \(W_t\succ 0\) et \(V_t\succ 0\) sont des poids de confiance, et \(\varepsilon_t,\eta_t>0\).
\end{definition}

\begin{lemma}[Convexité de la croyance]
L'ensemble \(\mathcal B_t\) est fermé, convexe et non vide dès que les contraintes observées sont compatibles avec \(\Omega_t\).
\end{lemma}

\begin{proof}
Chaque contrainte est l'inégalité d'une norme définissant une boule fermée convexe.  
L'intersection finie de fermés convexes est fermée et convexe.  
La non-vacuité suit de la compatibilité des observations avec l'état courant.
\end{proof}

\subsection{Énergie du monde}

\begin{definition}[Régularité du monde]
On pose
\[
\Omega_{\mathrm{world}}(\Omega)
=
\sum_{\ell=2}^{L}\|\omega^{(\ell)}-\omega^{(\ell-1)}\|^2.
\]
\end{definition}

\begin{definition}[Fonctionnelle dynamique du monde]
Soient \(o_t\in\Hh_{\mathrm{obs}}\) une observation et \(a_t\in\Hh_{\mathrm{act}}\) une action à l'instant \(t\).  
On définit
\[
J_{\mathrm{world}}(\Omega;\Omega_t,o_t,a_t)
=
\frac12\|\Omega-\Omega_t\|_{Q_t}^2
+
\frac12\|R\Omega-o_t\|_{W_t}^2
+
\frac12\|A\Omega-a_t\|_{V_t}^2
+
\lambda_{\mathrm{sm}}\Omega_{\mathrm{world}}(\Omega).
\]
\end{definition}

\begin{definition}[Mise à jour Multi-PH]
La mise à jour du monde est définie en deux étapes :
\[
\tilde\Omega_{t+1}
=
\argmin_{\Omega\in\Hh_\Omega}
J_{\mathrm{world}}(\Omega;\Omega_t,o_t,a_t),
\qquad
\Omega_{t+1}
=
\Pi_{\mathcal B_t}(\tilde\Omega_{t+1}).
\]
\end{definition}

\begin{proposition}[Existence et unicité de la mise à jour du monde]
Si \(Q_t\succ 0\), alors \(J_{\mathrm{world}}(\cdot;\Omega_t,o_t,a_t)\) est fortement convexe.  
Par conséquent, \(\tilde\Omega_{t+1}\) existe et est unique.  
Comme \(\mathcal B_t\) est fermé et convexe, la projection \(\Pi_{\mathcal B_t}\) est aussi unique, donc \(\Omega_{t+1}\) est unique.
\end{proposition}

\begin{proof}
Le terme \(\frac12\|\Omega-\Omega_t\|_{Q_t}^2\) est fortement convexe car \(Q_t\succ 0\).  
La somme avec les autres termes convexes reste fortement convexe.  
La projection orthogonale sur un fermé convexe est unique.
\end{proof}

\begin{definition}[Mise à jour de la confiance]
On définit la précision de croyance suivante :
\[
Q_{t+1}^{-1}
=
Q_t^{-1}
+
R^\top W_t R
+
A^\top V_t A.
\]
\end{definition}

%%%%%%%%%%%%%%%%%%%%%%%%%%%%%%%%%%%%%%%%%%%%%%%%%%%%%%%%%%%%%%%%%
%\begin{remark}
%Cette mise à jour rend l'incertitude explicite : plus les observations sont informatives, %plus la précision augmente.  
%La théorie ne masque donc pas l'incertitude ; elle la porte dans la croyance elle-même.
%\end{remark}
%%%%%%%%%%%%%%%%%%%%%%%%%%%%%%%%%%%%%%%%%%%%%%%%%%%%%%%%%%%%%%%%%

\begin{definition}[Itération proximale du monde]
Pour \(\eta>0\), on définit
\[
\mathcal T^\Omega_\eta(\Xi;\Omega_t,o_t,a_t)
=
\argmin_{\Omega\in\Hh_\Omega}
\left\{
J_{\mathrm{world}}(\Omega;\Omega_t,o_t,a_t)
+
\frac{1}{2\eta}\|\Omega-\Xi\|^2
\right\}.
\]
\end{definition}

\begin{proposition}[Contraction de la dynamique du monde]
Si \(J_{\mathrm{world}}(\cdot;\Omega_t,o_t,a_t)\) est \(\mu_{\mathrm{world}}\)-fortement convexe avec \(\mu_{\mathrm{world}}>0\), alors \(\mathcal T^\Omega_\eta(\cdot;\Omega_t,o_t,a_t)\) est une contraction.  
La trajectoire proximale converge vers l'unique point fixe, qui coïncide avec \(\tilde\Omega_{t+1}\).
\end{proposition}

\begin{proof}
La preuve est identique à celle de PH et PCH.  
La forte convexité donne la forte monotonie du sous-différentiel, ce qui entraîne la contraction du proximal.  
Banach donne alors la convergence vers l'unique point fixe.
\end{proof}

\begin{remark}
Dans cette perspective, prédire le prochain token n'est qu'un readout particulier d'un état latent hiérarchique.  
La théorie décrit donc un estimateur d'état certifié, et non une causalité cachée supplémentaire.
\end{remark}

\begin{definition}[Readout token]
Le prochain token est obtenu par une tête de lecture affine
\[
p_\psi(x_{t+1}\mid \Omega_{t+1})
=
\softmax(H_\psi \Omega_{t+1}),
\]
où \(H_\psi\) est une matrice apprise.
\end{definition}

% =================================================
\section{Fonction objective globale}
% =================================================

L'ensemble du système peut être entraîné ou calibré par une perte globale qui agrège les composantes validées.

\begin{definition}[Perte globale]
On définit
\[
\mathcal L_{\mathrm{tot}}
=
\lambda_{\mathrm{PCH}}\mathcal L_{\mathrm{PCH}}
+
\lambda_{\mathrm{spec}}\mathcal J_{\mathrm{spec}}
+
\lambda_{\mathrm{coord}}\mathcal J_{\mathrm{G}}
+
\lambda_{\mathrm{world}}\mathcal L_{\mathrm{world}}
+
\lambda_{\mathrm{task}}\mathcal L_{\mathrm{task}}
+
\lambda_{\mathrm{reg}}\Omega_{\mathrm{tot}}(\Theta,\Psi),
\]
où
\[
\mathcal L_{\mathrm{PCH}}
=
\mathbb E_{(x,c)}\!\left[
J_{\mathrm{PCH}}\bigl(\Pcal^{\mathrm{PCH}}(w(x);c);w(x),c\bigr)
\right],
\]
\[
\mathcal L_{\mathrm{world}}
=
\mathbb E_{t}\!\left[
J_{\mathrm{world}}(\Omega_{t+1};\Omega_t,o_t,a_t)
\right],
\]
et
\[
\mathcal J_{\mathrm{G}}
=
\mathbb E_{(x,c)}\!\left[
\mathcal J_{\mathrm{G}}\bigl(\bar u_0(x,c),\bar u_1(x,c),\dots,\bar u_n(x,c);x,c\bigr)
\right].
\]
\end{definition}

\begin{remark}
La perte globale ne crée pas de nouvelle boîte noire : elle agrège seulement les objectifs déjà explicitement définis.  
La partie essentielle de la coordination est le collage pondéré \(G\), pas une distillation cachée.
\end{remark}

% =================================================
\section{Résultats structurels}
% =================================================

\begin{proposition}[Principe d'unification structurelle]
Sous les hypothèses de convexité, de décomposition orthogonale spectrale, de séparabilité des domaines et de départage déterministe, l'architecture PCH-ESM est entièrement décrite par :
\begin{enumerate}[leftmargin=2em]
    \item un problème variationnel PH ;
    \item son extension contextuelle PCH ;
    \item une décomposition orthogonale du latent en sous-espaces spécialisés ;
    \item un objectif de spécialisation certifiée \(\mathcal J_{\mathrm{spec}}\) ;
    \item un collage pondéré certifié \(G\) sur la composante partagée ;
    \item une dynamique Multi-PH définie par minimisation d'une énergie explicite avec croyance et incertitude.
\end{enumerate}
\end{proposition}

\begin{proof}
Chaque composante est donnée explicitement : PH et PCH par des minimisations convexes, les sous-espaces par une construction spectrale, la spécialisation par une fonctionnelle définie sur un ensemble compact, \(G\) par un collage pondéré à retour, et Multi-PH par une minimisation énergétique suivie d'une projection sur un ensemble de croyance.  
Aucun objet structurel non défini n'est introduit.
\end{proof}

\begin{proposition}[Déterminisme global]
Si les ensembles de calibration, les partitions spectrales, les poids de confiance, les départages des minimisations et les paramètres sont fixés, alors la pipeline complète
\[
(x,c)\longmapsto \bigl(\hat y,\Omega_{t+1}\bigr)
\]
est déterministe.
\end{proposition}

\begin{proof}
Toutes les opérations intervenant dans la pipeline sont déterministes : moyenne, covariance, projecteurs spectraux, projections orthogonales, minimisations convexes uniques, collage pondéré, projection sur un fermé convexe, et lecture affine finale.  
La composition de ces opérations est donc déterministe.
\end{proof}

\begin{proposition}[Erreur globale contrôlée par les erreurs locales]
Soit
\[
\hat z=\bar u_0+\sum_{i=1}^n \bar u_i
\qquad\text{et}\qquad
z=z_0+\sum_{i=1}^n z_i,
\]
avec \(z_0,\bar u_0\in\Hh_0\) et \(z_i,\bar u_i\in\Hh_i\).  
Alors
\[
\|\hat z-z\|^2
=
\|\bar u_0-z_0\|^2+\sum_{i=1}^n\|\bar u_i-z_i\|^2.
\]
\end{proposition}

\begin{proof}
C'est la formule de Pythagore dans la somme directe orthogonale \(\Hh_0\oplus\cdots\oplus\Hh_n\).
\end{proof}

\begin{corollary}[Contrôle global]
Si
\[
\|\bar u_0-z_0\|\le \varepsilon_0
\quad\text{et}\quad
\|\bar u_i-z_i\|\le \varepsilon_i,
\]
alors
\[
\|\hat z-z\|
\le
\Bigl(\varepsilon_0^2+\sum_{i=1}^n \varepsilon_i^2\Bigr)^{1/2}.
\]
\end{corollary}

\begin{proof}
C'est une conséquence immédiate de la proposition précédente.
\end{proof}

\subsection{Capacité, dimension et nombre d'espaces}

\begin{definition}[Coût de coordination]
On définit la coordination moyenne associée à une configuration \((K,\mathbf d)\) par
\[
\mathcal E_{\mathrm{coord}}(K,\mathbf d)
=
\mathbb E_{(x,c)}\!\left[
\sum_{i=1}^{K}\bigl(1-\rho_i(x,c)\bigr)\,
\bigl\|\Psi_i(x,c)-\Pi_i E_{\theta_0}(x)\bigr\|^2
\right].
\]
\end{definition}

\begin{definition}[Fonction de capacité]
Soient \(K\in\{1,\dots,dT\}\) le nombre d'espaces spécialisés, \(\mathbf d=(d_1,\dots,d_K)\) une allocation avec \(\sum_i d_i=dT\), et \(\varepsilon_i:\mathbb N\to\R_+\) des erreurs locales décroissantes en dimension.  
On définit la fonction de capacité
\[
\mathcal C(K,\mathbf d)
=
\sum_{i=1}^K \varepsilon_i(d_i)
+
\lambda_{\mathrm{glue}}\mathcal E_{\mathrm{glue}}(K,\mathbf d)
+
\lambda_{\mathrm{coord}}\mathcal E_{\mathrm{coord}}(K,\mathbf d).
\]
\end{definition}

\begin{theorem}[Capacité et compromis entre dimension, nombre d'espaces et consensus]
La fonction \(\mathcal C(K,\mathbf d)\) atteint un minimum sur l'ensemble fini des allocations admissibles.  
De plus, le nombre optimal d'espaces \(K^\star\) existe, et il résulte d'un compromis explicite entre réduction de l'erreur locale et coût de consensus.
\end{theorem}

\begin{proof}
L'ensemble des allocations entières \(\mathbf d\) telles que \(\sum_i d_i=dT\) est fini pour un \(K\) fixé, et l'ensemble des valeurs possibles de \(K\) est fini.  
Ainsi, l'ensemble admissible total est fini, donc \(\mathcal C\) atteint un minimum.  
Le compromis local/consensus est donc non seulement conceptuel, mais aussi formellement défini par un problème de minimisation discret.
\end{proof}



% =================================================
\section{Interprétation théorique}
% =================================================

La théorie PCH-ESM affirme les points suivants :

\begin{enumerate}[leftmargin=2em]
    \item La bonne représentation n'est pas seulement celle qui prédit, mais celle qui reste stable sous purification récursive dans un contexte donné.
    \item L'attention standard peut être comprise comme un cas plat d'un schéma de sélection, sans qu'il soit nécessaire d'en faire une équivalence forte non démontrée.
    \item Les poids eux-mêmes peuvent être traités comme mémoire purifiable.
    \item Un espace sémantique monolithique peut être remplacé par une hiérarchie d'espaces spécialisés plus un noyau partagé.
    \item La généralisation globale ne disparaît pas nécessairement avec la spécialisation locale ; elle peut être reconstruite par un collage pondéré certifié sur la composante partagée.
    \item Les trajectoires internes de calcul deviennent explicites, donc traçables et déterministes dès lors que les départages sont fixés.
    \item La version monde du formalisme (Multi-PH) décrit un état hiérarchique estimé par minimisation variationnelle, avec croyance et incertitude.
    \item La capacité du système résulte d'un compromis formel entre dimension locale, nombre d'espaces et coût de consensus.
\end{enumerate}

\begin{remark}
Cette théorie ne suppose pas que la spécialisation soit gratuite : elle exige une couche de collage pour préserver la cohérence globale.
C'est précisément le rôle de \(G\).
\end{remark}

\section{Conclusion partielle}

PCH-ESM fournit une formalisation unifiée de :
\begin{itemize}[leftmargin=2em]
    \item la purification contextuelle hiérarchique ;
    \item la hiérarchisation de l'espace sémantique ;
    \item la décentralisation de l'IA en espaces spécialisés ;
    \item le collage pondéré certifié de la composante partagée ;
    \item le contrôle exact des erreurs globales ;
    \item la modélisation explicite de l'incertitude et des croyances ;
    \item le modèle du monde hiérarchique ;
    \item le compromis de capacité entre dimension, nombre d'espaces et consensus.
\end{itemize}

La proposition centrale est que la structure cognitive utile n'est pas uniquement un grand espace latent global, mais un \emph{écosystème hiérarchique d'espaces sémantiques} dont la stabilité, la spécialisation et la coordination sont explicitement représentées.





%%%%%%%%%%%%%%%%%%%%%%%%%%%%%%%%%%%%%%%%%%%%%%%%%%%%%%%%%%%%%%%%%%%%%%%%%%%%%%%%%%%%%%%%%%%%%%%%%%%%%%%%%%%%%%%%%%%%%%%%%%%%%%%%%%%%%%%%%%%%%%%%%%%%%%%%%%%%%%%%%%%%%%%%%%%%%%%%%%%%%%%%%%%%%%%%%%%%%%%%%%%%%%%%%%%%%%%%%%%%%%%%%%%%%%%%%%%%%%%%%%%%%%%%%%%%%%%%%%%%%%%%%%%%%%%%%%%%%%%%%%
%extention 1



% =================================================
\section{Résultats spécifiques à PCH-ESM}
% =================================================

Cette section établit quatre résultats qui ne sont pas de simples conséquences de la théorie proximale classique :
elle montre que la spécialisation réduit l'erreur sous compatibilité orthogonale,
que le collage pondéré surpasse la fusion naïve sous bruit hétéroscédastique,
que la partition spectrale canonique minimise exactement le couplage inter-blocs,
et que Multi-PH domine l'estimateur standard sous un modèle génératif cohérent avec la croyance hiérarchique.

% =================================================
\subsection{Gain de spécialisation certifié}
% =================================================

\begin{assumption}[Découplage orthogonal du problème d'apprentissage]
On suppose qu'il existe un sous-espace fermé \(\mathfrak S\) du paramètre global \(\Theta\) tel que
\[
\Theta=\mathfrak S\oplus \mathfrak S^\perp,
\]
et que l'objectif d'apprentissage quadratique pertinent pour la spécialisation se décompose orthogonalement :
\[
\mathcal J(\theta)
=
\mathcal J\bigl(\Pi_{\mathfrak S}\theta\bigr)
+
\mathcal J\bigl(\Pi_{\mathfrak S^\perp}\theta\bigr),
\]
où \(\Pi_{\mathfrak S}\) désigne la projection orthogonale sur \(\mathfrak S\).
On suppose en outre que le paramètre vrai \(\theta^\star\) appartient à \(\mathfrak S\).
\end{assumption}

\begin{theorem}[La spécialisation réduit l'erreur sous compatibilité orthogonale]
Soit \(\hat\theta_{\mathrm{dense}}\) le minimiseur global de \(\mathcal J\) dans \(\Theta\), et soit
\[
\hat\theta_{\mathrm{spec}}=\Pi_{\mathfrak S}\hat\theta_{\mathrm{dense}}
\]
le paramètre spécialisé obtenu par projection orthogonale.
Alors
\[
\|\hat\theta_{\mathrm{spec}}-\theta^\star\|
\le
\|\hat\theta_{\mathrm{dense}}-\theta^\star\|.
\]
De plus, si \(\Pi_{\mathfrak S^\perp}\hat\theta_{\mathrm{dense}}\neq 0\), alors l'inégalité est stricte.
\end{theorem}

\begin{proof}
Comme \(\theta^\star\in\mathfrak S\), on a la décomposition orthogonale
\[
\hat\theta_{\mathrm{dense}}-\theta^\star
=
\bigl(\Pi_{\mathfrak S}\hat\theta_{\mathrm{dense}}-\theta^\star\bigr)
+
\Pi_{\mathfrak S^\perp}\hat\theta_{\mathrm{dense}}.
\]
Les deux termes sont orthogonaux. Par Pythagore,
\[
\|\hat\theta_{\mathrm{dense}}-\theta^\star\|^2
=
\|\Pi_{\mathfrak S}\hat\theta_{\mathrm{dense}}-\theta^\star\|^2
+
\|\Pi_{\mathfrak S^\perp}\hat\theta_{\mathrm{dense}}\|^2.
\]
Or \(\hat\theta_{\mathrm{spec}}=\Pi_{\mathfrak S}\hat\theta_{\mathrm{dense}}\), donc
\[
\|\hat\theta_{\mathrm{dense}}-\theta^\star\|^2
=
\|\hat\theta_{\mathrm{spec}}-\theta^\star\|^2
+
\|\Pi_{\mathfrak S^\perp}\hat\theta_{\mathrm{dense}}\|^2.
\]
L'inégalité annoncée suit immédiatement, avec stricte inégalité si la composante orthogonale n'est pas nulle.
\end{proof}

\begin{corollary}[Interprétation PCH-ESM]
Sous les hypothèses du théorème précédent, la spécialisation ne constitue pas seulement une factorisation descriptive :
elle améliore effectivement la reconstruction du paramètre vrai dès lors que la solution dense contient une composante hors sous-espace spécialisé.
\end{corollary}

% =================================================
\subsection{Optimalité du collage pondéré}
% =================================================

\begin{assumption}[Modèle de bruit hétéroscédastique]
On suppose que les certificats locaux de la composante partagée s'écrivent
\[
u_0^{(i)} = u_0^\star + \varepsilon_i,
\qquad i=1,\dots,n,
\]
où \(u_0^\star\in\Hh_0\) est la vraie composante partagée, les erreurs \(\varepsilon_i\) sont indépendantes,
centrées, isotropes, et vérifient
\[
\mathbb E[\varepsilon_i]=0,
\qquad
\mathbb E\|\varepsilon_i\|^2 = d\,\tau_i,
\]
pour une variance scalaire \(\tau_i>0\).
\end{assumption}

\begin{theorem}[Le collage pondéré minimise l'erreur quadratique moyenne parmi les combinaisons affines non biaisées]
Considérons tout estimateur affine non biaisé de la forme
\[
\tilde u_0=\sum_{i=1}^n \beta_i u_0^{(i)},
\qquad
\sum_{i=1}^n \beta_i=1.
\]
Alors son erreur quadratique moyenne vaut
\[
\mathbb E\|\tilde u_0-u_0^\star\|^2
=
d\sum_{i=1}^n \beta_i^2 \tau_i.
\]
Cette quantité est minimisée lorsque
\[
\beta_i^\star
=
\frac{\tau_i^{-1}}{\sum_{j=1}^n \tau_j^{-1}},
\qquad i=1,\dots,n.
\]
Autrement dit, le consensus pondéré par des poids proportionnels aux fiabilités
\(\rho_i\propto \tau_i^{-1}\) est l'estimateur linéaire non biaisé optimal.
\end{theorem}

\begin{proof}
Par non biaiséness,
\[
\mathbb E[\tilde u_0-u_0^\star]
=
\sum_{i=1}^n \beta_i \mathbb E[\varepsilon_i]
=
0.
\]
Par indépendance et isotropie,
\[
\mathbb E\|\tilde u_0-u_0^\star\|^2
=
\mathbb E\left\|\sum_{i=1}^n \beta_i\varepsilon_i\right\|^2
=
\sum_{i=1}^n \beta_i^2\,\mathbb E\|\varepsilon_i\|^2
=
d\sum_{i=1}^n \beta_i^2\tau_i.
\]
On minimise cette quantité sous la contrainte linéaire \(\sum_i \beta_i=1\) par multiplicateurs de Lagrange :
\[
\mathcal L(\beta,\lambda)
=
d\sum_{i=1}^n \beta_i^2\tau_i
+
\lambda\left(\sum_{i=1}^n \beta_i-1\right).
\]
La condition stationnaire donne
\[
2d\,\tau_i\beta_i+\lambda=0,
\]
donc \(\beta_i\propto \tau_i^{-1}\). La contrainte de somme égale à \(1\) impose la formule annoncée.
\end{proof}

\begin{corollary}[G domine la fusion naïve lorsque les fiabilités sont inégales]
Si \(\tau_1,\dots,\tau_n\) ne sont pas toutes égales, alors le collage pondéré est strictement meilleur que la moyenne naïve
\[
u_0^{\mathrm{naif}}=\frac1n\sum_{i=1}^n u_0^{(i)}.
\]
Plus précisément,
\[
\mathbb E\|u_0^{G}-u_0^\star\|^2
\le
\mathbb E\|u_0^{\mathrm{naif}}-u_0^\star\|^2,
\]
avec stricte inégalité dès que les \(\tau_i\) ne sont pas toutes égales.
\end{corollary}

\begin{proof}
Le membre de gauche est l'erreur quadratique moyenne du minimiseur pondéré issu du théorème précédent.
La moyenne naïve correspond au choix \(\beta_i=\frac1n\), d'où
\[
\mathbb E\|u_0^{\mathrm{naif}}-u_0^\star\|^2
=
d\,\frac{1}{n^2}\sum_{i=1}^n \tau_i.
\]
Par Cauchy-Schwarz,
\[
\left(\sum_{i=1}^n \tau_i\right)\left(\sum_{i=1}^n \tau_i^{-1}\right)\ge n^2,
\]
donc
\[
\frac{1}{\sum_i \tau_i^{-1}}
\le
\frac1{n^2}\sum_i \tau_i.
\]
Or l'erreur du poids optimal vaut \(d(\sum_i\tau_i^{-1})^{-1}\). On obtient l'inégalité annoncée, avec stricte inégalité si les \(\tau_i\) ne sont pas toutes égales.
\end{proof}

\begin{corollary}[Interprétation PCH-ESM]
Dans le formalisme PCH-ESM, les poids \(\rho_i\) ne sont pas décoratifs :
ils réalisent la règle de pondération inverse-variance qui rend \(G\) optimal au sens quadratique parmi les collages affines non biaisés.
\end{corollary}

% =================================================
\subsection{Optimalité de la partition spectrale}
% =================================================

\begin{definition}[Couplage inter-blocs de la covariance]
Soient \(P_0,\dots,P_n\) des projecteurs orthogonaux de somme \(\Id\), définissant une partition orthogonale
\[
\Hh^T=\mathrm{Im}(P_0)\oplus\cdots\oplus\mathrm{Im}(P_n).
\]
On définit le couplage inter-blocs de la covariance empirique \(\Sigma\) par
\[
\mathcal J_{\mathrm{leak}}(P_0,\dots,P_n)
=
\sum_{i\neq j}\|P_i\Sigma P_j\|_F^2.
\]
\end{definition}

\begin{theorem}[La partition spectrale canonique minimise exactement le couplage inter-blocs]
Soit \(\Sigma\) symétrique positive semi-définie.  
Si \(P_0,\dots,P_n\) sont les projecteurs spectraux associés à la décomposition spectrale canonique de \(\Sigma\),
alors
\[
\mathcal J_{\mathrm{leak}}(P_0,\dots,P_n)=0.
\]
Par conséquent, cette partition est globalement optimale pour le critère de découplage inter-blocs :
\[
\mathcal J_{\mathrm{leak}}(Q_0,\dots,Q_n)\ge 0
\]
pour toute autre partition orthogonale \((Q_0,\dots,Q_n)\), avec égalité si et seulement si chaque \(Q_i\) est une somme de projecteurs spectraux de \(\Sigma\).
\end{theorem}

\begin{proof}
Par le théorème spectral, \(\Sigma\) admet une décomposition
\[
\Sigma=\sum_{k=1}^m \lambda_k P_k,
\]
où les \(P_k\) sont des projecteurs orthogonaux sur les sous-espaces propres.  
Si l'on regroupe ces projecteurs en blocs \(P_0,\dots,P_n\), alors chaque \(P_i\) commute avec \(\Sigma\), donc pour \(i\neq j\),
\[
P_i\Sigma P_j
=
P_iP_j\Sigma
=
0.
\]
Ainsi \(\mathcal J_{\mathrm{leak}}(P_0,\dots,P_n)=0\).  
Comme \(\mathcal J_{\mathrm{leak}}\) est somme de normes au carré, elle est toujours \(\ge 0\), donc la valeur \(0\) est globale minimale.  
L'égalité impose \(P_i\Sigma P_j=0\) pour tout \(i\neq j\), ce qui signifie précisément que les sous-espaces sont invariants par \(\Sigma\), donc constitués de sommes de sous-espaces propres.
\end{proof}

\begin{corollary}[Déterminisme et optimalité de la partition spectrale]
La partition spectrale déterministe utilisée dans PCH-ESM n'est pas seulement canonique :
elle est exactement celle qui annule le couplage inter-blocs de la covariance de calibration.
Elle constitue donc le meilleur choix pour séparer les blocs au sens du critère de leakage.
\end{corollary}

% =================================================
\subsection{Supériorité de Multi-PH sur l'estimateur standard}
% =================================================

\begin{assumption}[Modèle génératif gaussien hiérarchique]
On suppose que l'état vrai \(\Omega_{t+1}^\star\) suit le prior gaussien
\[
\Omega_{t+1}^\star \mid \Omega_t \sim \mathcal N(\Omega_t,Q_t^{-1}),
\]
et que les observations vérifient
\[
o_t = R\Omega_{t+1}^\star + \xi_t,
\qquad
a_t = A\Omega_{t+1}^\star + \zeta_t,
\]
où \(\xi_t\sim\mathcal N(0,W_t^{-1})\), \(\zeta_t\sim\mathcal N(0,V_t^{-1})\), indépendamment.
On suppose en outre que \(\Omega_{t+1}^\star\in\mathcal B_t\) presque sûrement.
\end{assumption}

\begin{definition}[Estimateur standard]
On définit l'estimateur standard par
\[
\hat\Omega_{t+1}^{\mathrm{std}}
=
\argmin_{\Omega\in\Hh_\Omega}
\left\{
\frac12\|R\Omega-o_t\|_{W_t}^2
+
\frac12\|A\Omega-a_t\|_{V_t}^2
\right\}.
\]
\end{definition}

\begin{definition}[Estimateur Multi-PH]
On définit l'estimateur Multi-PH par
\[
\tilde\Omega_{t+1}
=
\argmin_{\Omega\in\Hh_\Omega}
J_{\mathrm{world}}(\Omega;\Omega_t,o_t,a_t),
\qquad
\Omega_{t+1}^{\mathrm{MPH}}
=
\Pi_{\mathcal B_t}(\tilde\Omega_{t+1}).
\]
\end{definition}

\begin{theorem}[Multi-PH domine l'estimateur standard en risque quadratique]
Sous le modèle génératif gaussien ci-dessus, \(\tilde\Omega_{t+1}\) est l'estimateur MAP et, comme le posterior est gaussien, il coïncide avec l'estimateur bayésien pour la perte quadratique.
Par conséquent,
\[
\mathbb E\!\left[\|\tilde\Omega_{t+1}-\Omega_{t+1}^\star\|^2 \,\middle|\, o_t,a_t\right]
\le
\mathbb E\!\left[\|\hat\Omega_{t+1}^{\mathrm{std}}-\Omega_{t+1}^\star\|^2 \,\middle|\, o_t,a_t\right].
\]
De plus, comme \(\Omega_{t+1}^\star\in\mathcal B_t\), la projection finale ne peut pas augmenter l'erreur :
\[
\|\Omega_{t+1}^{\mathrm{MPH}}-\Omega_{t+1}^\star\|
\le
\|\tilde\Omega_{t+1}-\Omega_{t+1}^\star\|.
\]
Ainsi, Multi-PH est au moins aussi bon que l'estimateur standard, et strictement meilleur dès que le prior \((\Omega_t,Q_t)\) apporte une information non triviale.
\end{theorem}

\begin{proof}
Le terme \(J_{\mathrm{world}}(\Omega;\Omega_t,o_t,a_t)\) est, à constante additive près, le logarithme négatif de la densité a posteriori de \(\Omega_{t+1}^\star\) sous le modèle gaussien supposé.
Le minimiseur \(\tilde\Omega_{t+1}\) est donc l'estimateur MAP.  
Dans le cas gaussien quadratique, le MAP coïncide avec le posterior mean, lequel minimise le risque quadratique bayésien.
Cela implique la première inégalité.

Pour la seconde, on utilise le fait standard suivant :
si \(C\) est un fermé convexe et si \(y\in C\), alors pour tout \(x\),
\[
\|\Pi_C(x)-y\|^2+\|x-\Pi_C(x)\|^2\le \|x-y\|^2.
\]
En particulier, avec \(C=\mathcal B_t\) et \(y=\Omega_{t+1}^\star\in\mathcal B_t\), on obtient
\[
\|\Pi_{\mathcal B_t}(\tilde\Omega_{t+1})-\Omega_{t+1}^\star\|
\le
\|\tilde\Omega_{t+1}-\Omega_{t+1}^\star\|.
\]
Comme \(\Omega_{t+1}^{\mathrm{MPH}}=\Pi_{\mathcal B_t}(\tilde\Omega_{t+1})\), la conclusion suit.

La stricte amélioration suit dès que le prior n'est pas plat et que les observations sont bruitées, car alors le posterior mean diffère de l'estimateur standard sans prior.
\end{proof}

\begin{corollary}[Interprétation PCH-ESM]
Multi-PH n'est pas seulement un lissage proximal supplémentaire :
dans le modèle génératif explicitement posé, il est l'estimateur bayésien quadratique optimal, puis sa projection sur l'ensemble de croyance ne peut qu'améliorer ou préserver la précision.
\end{corollary}



%%%%%%%%%%%%%%%%%%%%%%%%%%%%%%%%%%%%%%%%%%%%%%%%%%%%%%%%%%%%%%%%%%%%%%%%%%%%%%%%%%%%%%%%%%%%%%%%%%%%%%%%%%%%%%%%%%%%%%%%%%%%%%%%%%%%%%%%%%%%%%%%%%%%%%%%%%%%%%%%%%%%%%%%%%%%%%%%%%%%%%%%%%%%%%%%%%%%%%%%%%%%%%%%%%%%%%%%%%%%%%%%%%%%%%%%%%%%%%%%%%%%%%%%%%%%%%%%%%%%%%%%%%%%%%%%%%%%%%%%%%
%extention 2

% =================================================
\subsection{Décomposition orthogonale exacte de l'erreur de reconstruction}
% =================================================

Cette sous-section ferme le verrou logique restant : elle montre que la décomposition
\[
\mathcal E_{\mathrm{rec}}(\mathbf P,G)
=
\mathcal E_{\mathrm{coll}}(\mathbf P,G)
+
\mathcal J_{\mathrm{leak}}(\mathbf P)
\]
résulte exactement de la décomposition orthogonale, pour la partition \(\mathbf P\), du résidu
entre l'opérateur cible et son approximation collée bloc-diagonale.

\begin{definition}[Espace des opérateurs et produit de Frobenius]
On pose
\[
\mathfrak E=\mathcal L(\Hh^T),
\]
muni du produit scalaire de Frobenius
\[
\langle A,B\rangle_F=\tr(A^\top B),
\qquad
\|A\|_F^2=\langle A,A\rangle_F.
\]
\end{definition}

\begin{definition}[Projecteurs de blocs]
Pour toute partition orthogonale admissible \(\mathbf P=(P_0,\dots,P_n)\), on définit les projections
bloc-diagonale et hors-diagonale par
\[
\Pi_{\mathrm{coll}}^{\mathbf P}(X)
=
\sum_{i=0}^n P_iXP_i,
\qquad
\Pi_{\mathrm{leak}}^{\mathbf P}(X)
=
\sum_{i\neq j}P_iXP_j,
\qquad
X\in\mathfrak E.
\]
On note aussi
\[
\mathfrak E_{\mathrm{coll}}^{\mathbf P}:=\mathrm{Im}\bigl(\Pi_{\mathrm{coll}}^{\mathbf P}\bigr),
\qquad
\mathfrak E_{\mathrm{leak}}^{\mathbf P}:=\mathrm{Im}\bigl(\Pi_{\mathrm{leak}}^{\mathbf P}\bigr).
\]
\end{definition}

\begin{lemma}[Orthogonalité des blocs]
Les sous-espaces \(\mathfrak E_{ij}^{\mathbf P}:=P_i\mathfrak E P_j\) sont deux à deux orthogonaux pour le produit de Frobenius.
En particulier,
\[
\mathfrak E
=
\bigoplus_{i,j=0}^n \mathfrak E_{ij}^{\mathbf P}
=
\mathfrak E_{\mathrm{coll}}^{\mathbf P}\oplus \mathfrak E_{\mathrm{leak}}^{\mathbf P},
\]
et les projections \(\Pi_{\mathrm{coll}}^{\mathbf P}\) et \(\Pi_{\mathrm{leak}}^{\mathbf P}\) sont orthogonales.
\end{lemma}

\begin{proof}
Soient \(A=P_iXP_j\) et \(B=P_kYP_\ell\). Alors
\[
\langle A,B\rangle_F
=
\tr\!\bigl((P_iXP_j)^\top P_kYP_\ell\bigr)
=
\tr\!\bigl(P_jX^\top P_iP_kYP_\ell\bigr).
\]
Si \(i\neq k\), alors \(P_iP_k=0\), donc \(\langle A,B\rangle_F=0\).

Si \(i=k\) et \(j\neq \ell\), alors
\[
\tr\!\bigl(P_jX^\top P_iYP_\ell\bigr)
=
\tr\!\bigl(P_\ell P_jX^\top P_iY\bigr)=0,
\]
car \(P_\ell P_j=0\). Ainsi, deux blocs distincts sont orthogonaux.

Enfin, comme \(\sum_{i=0}^n P_i=\Id\), tout opérateur \(X\in\mathfrak E\) se décompose en
\[
X=\sum_{i,j=0}^n P_iXP_j,
\]
d'où la somme directe orthogonale annoncée.
\end{proof}

\begin{definition}[Résidu de reconstruction et approximation collée]
On note \(\widehat \Sigma_{\mathrm{coll}}(G)\in\mathfrak E_{\mathrm{coll}}^{\mathbf P}\) l'opérateur reconstruit par la fonction de collage \(G\).

On définit le résidu total de reconstruction par
\[
R_{\mathrm{rec}}(\mathbf P,G)
=
\Sigma-\widehat \Sigma_{\mathrm{coll}}(G).
\]
On définit sa composante collée et sa composante de leakage par
\[
R_{\mathrm{coll}}(\mathbf P,G)
:=
\Pi_{\mathrm{coll}}^{\mathbf P}\!\bigl(R_{\mathrm{rec}}(\mathbf P,G)\bigr),
\qquad
R_{\mathrm{leak}}(\mathbf P)
:=
\Pi_{\mathrm{leak}}^{\mathbf P}(\Sigma).
\]
\end{definition}

\begin{definition}[Erreurs globales]
On définit
\[
\mathcal E_{\mathrm{rec}}(\mathbf P,G)
=
\mathbb E\bigl[\|R_{\mathrm{rec}}(\mathbf P,G)\|_F^2\bigr],
\qquad
\mathcal E_{\mathrm{coll}}(\mathbf P,G)
=
\mathbb E\bigl[\|R_{\mathrm{coll}}(\mathbf P,G)\|_F^2\bigr].
\]
Quand \(\mathbf P\) est fixé, on abrège parfois \(\mathcal E_{\mathrm{coll}}(G)\).
Le terme \(\mathcal E_{\mathrm{coll}}(\mathbf P,G)\) est l'erreur portée par la composante bloc-diagonale reconstruite.

L'espérance est prise relativement à la distribution des données de calibration utilisées pour construire
\(\widehat \Sigma_{\mathrm{coll}}(G)\).
\end{definition}

\begin{definition}[Le terme de leakage]
On pose
\[
\mathcal J_{\mathrm{leak}}(\mathbf P)
=
\bigl\|\Pi_{\mathrm{leak}}^{\mathbf P}(\Sigma)\bigr\|_F^2
=
\sum_{i\neq j}\|P_i\Sigma P_j\|_F^2.
\]
Le terme de leakage dépend uniquement de la covariance de calibration \(\Sigma\), supposée fixée ; il est donc déterministe.
\end{definition}

\begin{assumption}[Compatibilité minimale du mécanisme de collage]
Pour toute partition admissible \(\mathbf P\), il existe une constante \(C\ge 0\), indépendante de \(\mathbf P\), et une fonction strictement croissante
\[
\phi_{\mathbf P}:[0,\infty)\to[0,\infty),
\]
telles que, pour tout collage admissible \(G\),
\[
\mathcal E_{\mathrm{coll}}(\mathbf P,G)
=
C
+
\phi_{\mathbf P}\!\left(\mathbb E\bigl[\|\bar u_0(w,c)-u_0^\star(w,c)\|^2\bigr]\right).
\]
\end{assumption}

\begin{proposition}[Principe de réduction du collage]
Sous l'hypothèse précédente, pour toute partition admissible \(\mathbf P\),
\[
\argmin_G \mathcal E_{\mathrm{coll}}(\mathbf P,G)
=
\argmin_G \mathbb E\bigl[\|\bar u_0(w,c)-u_0^\star(w,c)\|^2\bigr].
\]
\end{proposition}

\begin{proof}
Comme \(\phi_{\mathbf P}\) est strictement croissante, minimiser \(\mathcal E_{\mathrm{coll}}(\mathbf P,G)\) équivaut exactement à minimiser
\[
\mathbb E\bigl[\|\bar u_0(w,c)-u_0^\star(w,c)\|^2\bigr].
\]
La constante \(C\) est indépendante de \(G\) et de \(\mathbf P\), donc elle n'affecte pas l'ensemble des minimiseurs.
\end{proof}

\begin{theorem}[Décomposition orthogonale exacte de l'erreur de reconstruction]
Sous les définitions précédentes, pour toute architecture admissible \((\mathbf P,G)\),
\[
\mathcal E_{\mathrm{rec}}(\mathbf P,G)
=
\mathcal E_{\mathrm{coll}}(\mathbf P,G)
+
\mathcal J_{\mathrm{leak}}(\mathbf P).
\]
\end{theorem}

\begin{proof}
On écrit d'abord
\[
R_{\mathrm{rec}}(\mathbf P,G)
=
\Sigma-\widehat \Sigma_{\mathrm{coll}}(G).
\]
Comme \(\widehat \Sigma_{\mathrm{coll}}(G)\in \mathfrak E_{\mathrm{coll}}^{\mathbf P}\), sa projection hors-diagonale est nulle, donc
\[
\Pi_{\mathrm{leak}}^{\mathbf P}\!\bigl(R_{\mathrm{rec}}(\mathbf P,G)\bigr)
=
\Pi_{\mathrm{leak}}^{\mathbf P}(\Sigma).
\]
De plus,
\[
R_{\mathrm{rec}}(\mathbf P,G)
=
\underbrace{\Bigl(\Pi_{\mathrm{coll}}^{\mathbf P}(\Sigma)-\widehat \Sigma_{\mathrm{coll}}(G)\Bigr)}_{R_{\mathrm{coll}}(\mathbf P,G)\in\mathfrak E_{\mathrm{coll}}^{\mathbf P}}
+
\underbrace{\Pi_{\mathrm{leak}}^{\mathbf P}(\Sigma)}_{R_{\mathrm{leak}}(\mathbf P)\in\mathfrak E_{\mathrm{leak}}^{\mathbf P}}.
\]
Les deux termes appartiennent à des sous-espaces orthogonaux, donc le terme croisé s'annule exactement :
\[
\|R_{\mathrm{rec}}(\mathbf P,G)\|_F^2
=
\|R_{\mathrm{coll}}(\mathbf P,G)\|_F^2
+
\|R_{\mathrm{leak}}(\mathbf P)\|_F^2.
\]
En prenant l'espérance,
\[
\mathcal E_{\mathrm{rec}}(\mathbf P,G)
=
\mathbb E\bigl[\|R_{\mathrm{coll}}(\mathbf P,G)\|_F^2\bigr]
+
\|R_{\mathrm{leak}}(\mathbf P)\|_F^2.
\]
Par définition,
\[
\mathcal E_{\mathrm{coll}}(\mathbf P,G)
=
\mathbb E\bigl[\|R_{\mathrm{coll}}(\mathbf P,G)\|_F^2\bigr].
\]
En outre,
\[
\|R_{\mathrm{leak}}(\mathbf P)\|_F^2
=
\bigl\|\Pi_{\mathrm{leak}}^{\mathbf P}(\Sigma)\bigr\|_F^2
=
\sum_{i\neq j}\|P_i\Sigma P_j\|_F^2
=
\mathcal J_{\mathrm{leak}}(\mathbf P).
\]
Donc finalement,
\[
\mathcal E_{\mathrm{rec}}(\mathbf P,G)
=
\mathcal E_{\mathrm{coll}}(\mathbf P,G)
+
\mathcal J_{\mathrm{leak}}(\mathbf P),
\]
ce qu'il fallait démontrer.
\end{proof}

%%%%%%%%%%%%%%%%%%%%%%%%%%%%%%%%%%%%%%%%%%%%%%%%%%%%%%%%%%%%%%%%%%%%
%\begin{remark}
%La décomposition orthogonale précédente est indépendante du mécanisme de collage.
%En revanche, l'identification du minimiseur de \(\mathcal E_{\mathrm{coll}}\) avec le collage %inverse-variance repose uniquement sur l'hypothèse minimale de compatibilité du mécanisme de %collage, qui exprime que \(\mathcal E_{\mathrm{coll}}\) est une transformation strictement %croissante de l'erreur de partage.

%Le théorème précédent remplace l'hypothèse de décomposition leakage--reconstruction utilisée %dans les versions préliminaires du formalisme.

%\end{remark}
%%%%%%%%%%%%%%%%%%%%%%%%%%%%%%%%%%%%%%%%%%%%%%%%%%%%%%%%%%%%%%%%%%%%


% =================================================
\subsection{Optimalité conjointe sous compatibilité spectrale exacte}
% =================================================

Cette sous-section établit le résultat structural maximal du formalisme sous une hypothèse de compatibilité exacte entre
(i) le découplage spectral des modules,
(ii) l'ordre hiérarchique des blocs, et
(iii) la reconstruction globale par collage pondéré.

\begin{assumption}[Spectre simple et partitions orthogonales de rangs fixés]
On suppose que la covariance de calibration \(\Sigma\) est symétrique positive semi-définie
et que son spectre est simple.
On ordonne ses valeurs propres de façon décroissante :
\[
\lambda_1>\lambda_2>\cdots>\lambda_{dT}\ge 0,
\]
avec une base orthonormée associée \((e_1,\dots,e_{dT})\) telle que
\[
\Sigma e_k=\lambda_k e_k.
\]

On fixe des tailles entières \(d_0,\dots,d_n\ge 0\) telles que
\[
d_0+\cdots+d_n=dT.
\]
On note
\[
I_0=\{1,\dots,d_0\},\quad
I_1=\{d_0+1,\dots,d_0+d_1\},\quad \dots,\quad
I_n=\{d_0+\cdots+d_{n-1}+1,\dots,dT\}.
\]

On définit la partition canonique
\[
\mathbf P^\star=(P_0^\star,\dots,P_n^\star)
\]
par
\[
P_i^\star=\sum_{k\in I_i} e_ke^\top_k,
\qquad i=0,\dots,n.
\]

On appelle \emph{admissible} toute partition orthogonale
\[
\mathbf P=(P_0,\dots,P_n)
\]
de rangs fixés vérifiant
\[
P_iP_j=0\quad(i\neq j),
\qquad
\sum_{i=0}^n P_i=\Id.
\]
La classe admissible n'est donc pas réduite a priori à \(\{\mathbf P^\star\}\) :
l'unicité sera obtenue par le critère d'optimalité ci-dessous.
\end{assumption}

\begin{assumption}[Compatibilité hiérarchique exacte]
Les architectures admissibles sont les couples \((\mathbf P,G)\) tels que :

\begin{itemize}[leftmargin=2em]
    \item \(\mathbf P=(P_0,\dots,P_n)\) est une partition orthogonale admissible ;
    \item les certificats locaux appartiennent aux sous-espaces associés, c'est-à-dire
    \[
    u_0^{(i)}(w,c)\in \mathrm{Im}(P_0),
    \qquad
    \Psi_i(w,c)\in \mathrm{Im}(P_i);
    \]
    \item le collage global \(G\) est affine en les certificats locaux, et son estimateur partagé
    \(\bar u_0(w,c)\) est non biaisé au sens
    \[
    \mathbb E[\bar u_0(w,c)]=u_0^\star(w,c).
    \]
\end{itemize}
\end{assumption}

\begin{assumption}[Modèle de bruit hétéroscédastique]
Pour chaque \(i\in\{1,\dots,n\}\), les certificats partagés vérifient
\[
u_0^{(i)}=u_0^\star+\varepsilon_i,
\]
où les \(\varepsilon_i\) sont indépendantes, centrées, isotropes, et de variance scalaire
\[
\mathbb E\|\varepsilon_i\|^2=d\,\tau_i,
\qquad \tau_i>0.
\]
\end{assumption}

\begin{assumption}[Poids strictement positifs]
On suppose
\[
\lambda_{\mathrm{leak}}>0,
\qquad
\lambda_{\mathrm{ord}}>0,
\qquad
\lambda_{\mathrm{rec}}>0.
\]
\end{assumption}

\begin{definition}[Terme d'ordre spectral]
On fixe des poids hiérarchiques strictement décroissants
\[
\omega_0>\omega_1>\cdots>\omega_n>0.
\]
On définit le terme d'ordre par
\[
\mathcal J_{\mathrm{ord}}(\mathbf P)
=
-\sum_{i=0}^n \omega_i\,\mathrm{tr}(P_i\Sigma).
\]
Ce terme impose que les plus grandes directions spectrales soient affectées aux blocs les plus prioritaires.
\end{definition}

\begin{definition}[Objectif architectural conjoint]
On définit l'objectif global par
\[
\mathcal J_{\mathrm{arch}}(\mathbf P,G)
=
\lambda_{\mathrm{leak}}\,\mathcal J_{\mathrm{leak}}(\mathbf P)
+
\lambda_{\mathrm{ord}}\,\mathcal J_{\mathrm{ord}}(\mathbf P)
+
\lambda_{\mathrm{rec}}\,\mathcal E_{\mathrm{rec}}(\mathbf P,G).
\]
En utilisant le théorème précédent, on obtient l'écriture équivalente
\[
\mathcal J_{\mathrm{arch}}(\mathbf P,G)
=
(\lambda_{\mathrm{leak}}+\lambda_{\mathrm{rec}})\,
\mathcal J_{\mathrm{leak}}(\mathbf P)
+
\lambda_{\mathrm{ord}}\,\mathcal J_{\mathrm{ord}}(\mathbf P)
+
\lambda_{\mathrm{rec}}\,\mathcal E_{\mathrm{coll}}(\mathbf P,G).
\]
\end{definition}

\begin{lemma}[Zéro leakage et compatibilité spectrale]
Pour toute partition admissible \(\mathbf P\),
\[
\mathcal J_{\mathrm{leak}}(\mathbf P)\ge 0.
\]
De plus,
\[
\mathcal J_{\mathrm{leak}}(\mathbf P)=0
\quad\Longleftrightarrow\quad
P_i\Sigma=\Sigma P_i
\ \text{ pour tout } i,
\]
c'est-à-dire si et seulement si chaque bloc \(P_i\) est une somme de projecteurs spectraux de \(\Sigma\).
\end{lemma}

\begin{proof}
Par définition, \(\mathcal J_{\mathrm{leak}}\) est une somme de normes au carré, donc elle est non négative.
Si elle vaut \(0\), alors \(P_i\Sigma P_j=0\) pour tous \(i\neq j\), ce qui signifie que \(\Sigma\) est bloc-diagonale dans la décomposition \(\mathbf P\), donc que chaque \(P_i\) commute avec \(\Sigma\).
Réciproquement, si chaque \(P_i\) commute avec \(\Sigma\), alors les blocs hors diagonale sont nuls et \(\mathcal J_{\mathrm{leak}}(\mathbf P)=0\).
\end{proof}

\begin{lemma}[Réécriture télescopique du terme d'ordre]
Soient
\[
m_r=\sum_{i=0}^r d_i
\qquad\text{et}\qquad
Q_r=\sum_{i=0}^r P_i,
\quad r=0,\dots,n,
\]
avec la convention \(Q_{-1}=0\).
Alors
\[
\sum_{i=0}^n \omega_i\,\mathrm{tr}(P_i\Sigma)
=
\omega_n\,\mathrm{tr}(\Sigma)
+
\sum_{r=0}^{n-1}(\omega_r-\omega_{r+1})\,\mathrm{tr}(Q_r\Sigma).
\]
\end{lemma}

\begin{proof}
Comme \(P_i=Q_i-Q_{i-1}\), on a
\[
\sum_{i=0}^n \omega_i\,\mathrm{tr}(P_i\Sigma)
=
\sum_{i=0}^n \omega_i\bigl(\mathrm{tr}(Q_i\Sigma)-\mathrm{tr}(Q_{i-1}\Sigma)\bigr).
\]
Un regroupement télescopique donne
\[
\sum_{i=0}^n \omega_i\,\mathrm{tr}(P_i\Sigma)
=
\omega_n\,\mathrm{tr}(Q_n\Sigma)
+
\sum_{r=0}^{n-1}(\omega_r-\omega_{r+1})\,\mathrm{tr}(Q_r\Sigma).
\]
Comme \(Q_n=\Id\), on obtient la formule annoncée.
\end{proof}

\begin{lemma}[Sélection canonique par l'ordre spectral]
Parmi les partitions admissibles, le terme
\[
\mathcal J_{\mathrm{ord}}(\mathbf P)
=
-\sum_{i=0}^n \omega_i\,\mathrm{tr}(P_i\Sigma)
\]
est strictement minimisé par la partition canonique \(\mathbf P^\star\), et seulement par elle.
\end{lemma}

\begin{proof}
Par le lemme précédent, minimiser \(\mathcal J_{\mathrm{ord}}\) revient à maximiser chaque quantité \(\mathrm{tr}(Q_r\Sigma)\), puisque les coefficients \(\omega_r-\omega_{r+1}\) sont strictement positifs.

Pour tout projecteur orthogonal \(Q\) de rang \(m_r\), le principe de Ky Fan donne
\[
\mathrm{tr}(Q\Sigma)\le \sum_{k=1}^{m_r}\lambda_k,
\]
avec égalité si et seulement si
\[
Q=\sum_{k=1}^{m_r} e_k e_k^\top,
\]
l'unicité étant assurée par la simplicité du spectre.

En l'appliquant à chaque \(Q_r\), on obtient
\[
\mathrm{tr}(Q_r\Sigma)\le \sum_{k=1}^{m_r}\lambda_k,
\]
et l'égalité simultanée pour tout \(r\) impose
\[
Q_r=\sum_{k=1}^{m_r} e_k e_k^\top
\qquad\text{pour tout } r.
\]
Par différence,
\[
P_i=Q_i-Q_{i-1}
=
\sum_{k\in I_i} e_ke^\top_k
=
P_i^\star.
\]
Donc \(\mathbf P^\star\) est l'unique minimiseur de \(\mathcal J_{\mathrm{ord}}\) sur la classe admissible.
\end{proof}

\begin{proposition}[Optimalité du collage inverse-variance]
Sous le modèle de bruit hétéroscédastique, tout estimateur affine non biaisé de la forme
\[
\tilde u_0=\sum_{i=1}^n \beta_i u_0^{(i)},
\qquad
\sum_{i=1}^n \beta_i=1,
\]
vérifie
\[
\mathbb E\|\tilde u_0-u_0^\star\|^2
=
d\sum_{i=1}^n \beta_i^2\tau_i.
\]
Cette quantité est strictement minimisée par
\[
\beta_i^\star
=
\frac{\tau_i^{-1}}{\sum_{j=1}^n\tau_j^{-1}},
\qquad i=1,\dots,n.
\]
\end{proposition}

\begin{proof}
On a
\[
\mathbb E[\tilde u_0-u_0^\star]
=
\sum_{i=1}^n \beta_i\mathbb E[\varepsilon_i]
=
0.
\]
Par indépendance et isotropie,
\[
\mathbb E\|\tilde u_0-u_0^\star\|^2
=
\mathbb E\left\|\sum_{i=1}^n\beta_i\varepsilon_i\right\|^2
=
\sum_{i=1}^n\beta_i^2\,\mathbb E\|\varepsilon_i\|^2
=
d\sum_{i=1}^n\beta_i^2\tau_i.
\]
La fonction
\[
\beta\mapsto d\sum_{i=1}^n\beta_i^2\tau_i
\]
est strictement convexe sur l'hyperplan \(\sum_i\beta_i=1\).  
Le minimiseur unique est donc obtenu par multiplicateurs de Lagrange, ce qui donne
\[
\beta_i^\star
=
\frac{\tau_i^{-1}}{\sum_{j=1}^n\tau_j^{-1}}.
\]
\end{proof}

\begin{theorem}[Optimalité conjointe et unicité du couple partition spectrale + collage pondéré]
Sous les hypothèses précédentes, le problème
\[
\min_{(\mathbf P,G)\ \mathrm{admissible}} \ \mathcal J_{\mathrm{arch}}(\mathbf P,G)
\]
admet un unique minimiseur admissible.
Ce minimiseur est précisément le couple
\[
(\mathbf P^\star,G^\star),
\]
où :

\begin{itemize}[leftmargin=2em]
    \item \(\mathbf P^\star\) est la partition spectrale canonique ;
    \item \(G^\star\) est le collage pondéré inverse-variance, i.e.
    \[
    \bar u_0^\star(w,c)
    =
    \sum_{i=1}^n \alpha_i^\star(w,c)\,u_0^{(i)}(w,c),
    \qquad
    \alpha_i^\star(w,c)
    =
    \frac{\tau_i^{-1}}{\sum_{j=1}^n \tau_j^{-1}}.
    \]
\end{itemize}

En particulier, le couple \((\mathbf P^\star,G^\star)\) est l'unique architecture admissible qui minimise simultanément le leakage inter-modules et l'erreur globale de reconstruction.
\end{theorem}

\begin{proof}
Posons
\[
A=\lambda_{\mathrm{leak}}+\lambda_{\mathrm{rec}}>0.
\]
Alors, pour tout couple admissible \((\mathbf P,G)\),
\[
\mathcal J_{\mathrm{arch}}(\mathbf P,G)
=
A\,\mathcal J_{\mathrm{leak}}(\mathbf P)
+
\lambda_{\mathrm{ord}}\,\mathcal J_{\mathrm{ord}}(\mathbf P)
+
\lambda_{\mathrm{rec}}\,\mathcal E_{\mathrm{coll}}(\mathbf P,G).
\]

\medskip
\noindent
\textbf{Étape 1 : optimisation de la partie partition.}
Comme \(\mathcal J_{\mathrm{leak}}(\mathbf P)\ge 0\) et que \(\mathbf P^\star\) minimise strictement \(\mathcal J_{\mathrm{ord}}(\mathbf P)\) sur la classe admissible, on a pour toute partition admissible
\[
A\,\mathcal J_{\mathrm{leak}}(\mathbf P)+\lambda_{\mathrm{ord}}\,\mathcal J_{\mathrm{ord}}(\mathbf P)
\ge
\lambda_{\mathrm{ord}}\,\mathcal J_{\mathrm{ord}}(\mathbf P)
\ge
\lambda_{\mathrm{ord}}\,\mathcal J_{\mathrm{ord}}(\mathbf P^\star).
\]
De plus,
\[
\mathcal J_{\mathrm{leak}}(\mathbf P^\star)=0
\]
par compatibilité spectrale de la partition canonique.
Ainsi, la partie dépendant de \(\mathbf P\) est minimisée de manière unique par \(\mathbf P^\star\).

\medskip
\noindent
\textbf{Étape 2 : optimisation du collage.}
Une fois \(\mathbf P^\star\) fixé, la partie dépendant de \(G\) se réduit à \(\mathcal E_{\mathrm{coll}}(\mathbf P^\star,G)\).
Par la proposition de réduction du collage,
\[
\argmin_G \mathcal E_{\mathrm{coll}}(\mathbf P^\star,G)
=
\argmin_G \mathbb E\bigl[\|\bar u_0(w,c)-u_0^\star(w,c)\|^2\bigr].
\]
Sous le modèle de bruit hétéroscédastique, la proposition précédente montre que cet ensemble est réduit au seul collage affine non biaisé inverse-variance \(G^\star\), avec
\[
\alpha_i^\star(w,c)=\frac{\tau_i^{-1}}{\sum_{j=1}^n\tau_j^{-1}}.
\]

\medskip
\noindent
\textbf{Étape 3 : unicité globale.}
Si \((\mathbf P,G)\) est un minimiseur global, alors il doit satisfaire simultanément
\[
\mathbf P=\mathbf P^\star
\qquad\text{et}\qquad
G=G^\star,
\]
car l'égalité dans l'étape 1 force \(\mathbf P=\mathbf P^\star\), et l'égalité dans l'étape 2 force \(G=G^\star\).
Réciproquement, \((\mathbf P^\star,G^\star)\) atteint cette borne minimale.
Le minimiseur global est donc unique et vaut exactement \((\mathbf P^\star,G^\star)\).
\end{proof}

%%%%%%%%%%%%%%%%%%%%%%%%%%%%%%%%%%%%%%%%%%%%%%%%%%%%%%%%%%%%%%%%%%%%
\begin{remark}
La décomposition orthogonale précédente est indépendante du mécanisme de collage.
En revanche, l'identification du minimiseur de \(\mathcal E_{\mathrm{coll}}\) avec le collage inverse-variance repose uniquement sur l'hypothèse minimale de compatibilité du mécanisme de collage, qui exprime que \(\mathcal E_{\mathrm{coll}}\) est une transformation strictement croissante de l'erreur de partage.

%Le théorème précédent remplace l'hypothèse de décomposition leakage--reconstruction utilisée %dans les versions préliminaires du formalisme.

\end{remark}
%%%%%%%%%%%%%%%%%%%%%%%%%%%%%%%%%%%%%%%%%%%%%%%%%%%%%%%%%%%%%%%%%%%%

\begin{corollary}[Optimalité conjointe]
Le formalisme PCH-ESM réalise simultanément :
\begin{enumerate}[leftmargin=2em]
    \item le découplage maximal des modules via la partition spectrale canonique ;
    \item la reconstruction globale de risque minimal via le collage pondéré inverse-variance.
\end{enumerate}
Autrement dit, sous les hypothèses précédentes, le couple \((\mathbf P^\star,G^\star)\) est l'unique minimiseur simultané du leakage inter-modules et de l'erreur de reconstruction dans la classe des architectures admissibles.
\end{corollary}


% =================================================
\section{Portée et limites}
% =================================================

Cette théorie établit une cohérence interne forte et donne des résultats exacts sur :
\begin{itemize}[leftmargin=2em]
    \item l'existence et l'unicité des purifications PH, PCH et Multi-PH ;
    \item la convergence des trajectoires proximales ;
    \item la construction explicite des espaces spécialisés ;
    \item le collage pondéré certifié de la composante partagée ;
    \item le contrôle exact de l'erreur globale par les erreurs locales ;
    \item la modélisation explicite de l'incertitude et des croyances ;
    \item l'existence d'un optimum de capacité pour le nombre d'espaces.
\end{itemize}

En revanche, cette formalisation ne prétend pas démontrer à elle seule une supériorité empirique universelle sur toutes les architectures existantes.  
Cette question relève ensuite de l'évaluation expérimentale.


\section{Conclusion}

PCH-ESM propose une lecture structurée de la représentation : la bonne architecture n’est pas seulement prédictive, elle est stable, spécialisée, coordonnée et explicitement réconciliée dans un cadre hiérarchique.

Le noyau théorique du document est maintenant établi : PH et PCH donnent des opérateurs de purification bien posés et contractants ; la décomposition spectrale fournit des sous-espaces spécialisés déterministes ; la spécialisation certifiée sélectionne les blocs pertinents ; le collage pondéré inverse-variance reconstruit la composante partagée optimale ; la décomposition orthogonale de l’erreur sépare proprement l’effet du collage et le leakage ; enfin, Multi-PH étend le formalisme à un état du monde hiérarchique avec croyance et incertitude.

Les résultats structurels montrent que, sous les hypothèses du modèle, la partition spectrale canonique et le collage pondéré optimal réalisent simultanément le découplage maximal des modules et la reconstruction de risque minimal. La théorie donne aussi un contrôle explicite du compromis entre dimension locale, nombre d’espaces et coût de consensus.

La portée de ces résultats est interne et formelle : elle établit une cohérence mathématique forte, mais ne remplace pas l’évaluation empirique. La théorie fournit donc une base structurale complète pour étudier expérimentalement quand et pourquoi une représentation hiérarchique spécialisée, pondérée et contextuelle peut être préférable à une architecture monolithique.




% =================================================
\appendix
\section{Validation expérimentale du noyau proximal PCH}
\label{app:pch-proximal-poc}
% =================================================

\subsection{Objet de l'expérience}

Cet appendice formalise la preuve de concept numérique du noyau proximal de PCH sur un espace de Hilbert fini-dimensionnel.
L'objectif n'est pas de valider l'ensemble complet de la théorie PCH-ESM, mais d'établir rigoureusement les points suivants :

\begin{itemize}
    \item l'existence d'un minimiseur unique pour le problème proximal PCH ;
    \item la forme fermée de l'opérateur proximal ;
    \item la convergence vers le point fixe exact ;
    \item la décroissance monotone de l'énergie ;
    \item une estimation correcte de la contraction locale du map proximal, mesurée sur les premières itérations utiles, avant saturation numérique.
\end{itemize}

L'expérience porte donc sur le noyau proximal exact
\[
u^\star(w,c)
=
\operatorname*{argmin}_{u\in\mathcal H}
\left\{
\frac12\|u-w\|^2 + \Phi_{\mathrm{PCH},c}(u)
\right\},
\]
ainsi que sur l'itération proximale associée
\[
u^{(r+1)} = T_\eta(u^{(r)};w,c),
\]
où \(T_\eta\) est calculé en forme fermée.

\subsection{Modèle mathématique exact}

On considère un espace de Hilbert fini-dimensionnel \(\mathcal H \simeq \mathbb R^d\), muni du produit scalaire euclidien.
Pour chaque contexte \(c\), on fixe un sous-espace fermé \(M_c\subset\mathcal H\), modélisé dans l'implémentation par un projecteur orthogonal \(P_c\), avec
\[
M_c = \operatorname{Im}(P_c),
\qquad
P_c^2=P_c,
\qquad
P_c^\top=P_c.
\]
Le fonctionnel PCH testé numériquement est
\[
\Phi_{\mathrm{PCH},c}(u)
=
\lambda_{\mathrm{ctx}}\,\operatorname{dist}(u,M_c)^2
+
\lambda_{\mathrm{ch}}\|u\|^2
+
\lambda_{\mathrm{reg}}\|u\|^2.
\]
Comme
\[
\operatorname{dist}(u,M_c)^2 = \|(I-P_c)u\|^2,
\]
on obtient un problème quadratique strictement convexe, de minimiseur unique.

En posant
\[
\lambda_{\mathrm{tot}}=\lambda_{\mathrm{ch}}+\lambda_{\mathrm{reg}},
\]
le problème statique
\[
u^\star(w,c)
=
\operatorname*{argmin}_{u}
\left\{
\frac12\|u-w\|^2+\Phi_{\mathrm{PCH},c}(u)
\right\}
\]
admet la solution fermée
\[
u^\star(w,c)
=
\frac{P_c w}{1+2\lambda_{\mathrm{tot}}}
+
\frac{(I-P_c)w}{1+2(\lambda_{\mathrm{ctx}}+\lambda_{\mathrm{tot}})}.
\]

L'opérateur proximal régularisé est défini par
\[
T_\eta(z;w,c)
=
\operatorname*{argmin}_{u}
\left\{
\frac12\|u-w\|^2+\Phi_{\mathrm{PCH},c}(u)+\frac{1}{2\eta}\|u-z\|^2
\right\}.
\]
Dans la décomposition orthogonale \(\mathcal H=M_c\oplus M_c^\perp\), ce map se sépare exactement en
\[
T_\eta(z;w,c)
=
\frac{P_c w + \eta^{-1}P_c z}{1+2\lambda_{\mathrm{tot}}+\eta^{-1}}
+
\frac{(I-P_c)w + \eta^{-1}(I-P_c)z}{1+2(\lambda_{\mathrm{ctx}}+\lambda_{\mathrm{tot}})+\eta^{-1}}.
\]

Le point fixe exact de l'itération est \(u^\star(w,c)\), et le noyau linéarisé en \(z\) est diagonal sur la décomposition \(M_c\oplus M_c^\perp\). Sa norme opérateur exacte vaut
\[
\rho(\eta)
=
\max\left\{
\frac{1}{1+\eta(1+2\lambda_{\mathrm{tot}})},
\frac{1}{1+\eta(1+2(\lambda_{\mathrm{ctx}}+\lambda_{\mathrm{tot}}))}
\right\}.
\]
Cette quantité est la vraie constante de contraction à comparer expérimentalement, mais uniquement sur les premières itérations utiles, avant que l'erreur ne tombe au voisinage du bruit machine.

\subsection{Protocole expérimental}

Le protocole suivi est le suivant.

\begin{enumerate}
    \item Construire de manière déterministe des sous-espaces contextuels \(M_c\) via des bases orthonormées générées à partir d'une graine fixée.
    \item Fixer une famille de contextes \(c\in\{c_1,\dots,c_K\}\).
    \item Échantillonner des vecteurs \(w\) d'entrée.
    \item Calculer, pour chaque paire \((w,c)\), le minimiseur exact \(u^\star(w,c)\).
    \item Initialiser l'itération par \(u^{(0)}=w\), puis appliquer
    \[
    u^{(r+1)}=T_\eta(u^{(r)};w,c).
    \]
    \item Enregistrer à chaque étape l'énergie
    \[
    J_{\mathrm{PCH}}(u^{(r)};w,c)
    =
    \frac12\|u^{(r)}-w\|^2+\Phi_{\mathrm{PCH},c}(u^{(r)}),
    \]
    l'erreur au point fixe \(\|u^{(r)}-u^\star\|\), et le résidu \(\|T_\eta(u^\star;w,c)-u^\star\|\).
    \item Calculer les ratios de contraction successifs
    \[
    \rho_r
    =
    \frac{\|u^{(r+1)}-u^\star\|}{\|u^{(r)}-u^\star\|},
    \]
    mais uniquement pour les indices \(r\) tels que
    \[
    \|u^{(r)}-u^\star\| \ge \tau_{\mathrm{fp}},
    \]
    où \(\tau_{\mathrm{fp}}\) est un seuil de filtrage strictement au-dessus du bruit machine.
\end{enumerate}

Le point important est le suivant : la mesure de contraction ne doit pas être moyennée sur toute la trajectoire sans filtre, car les derniers termes sont dominés par la précision flottante et ne mesurent plus la dynamique théorique. La bonne comparaison empirique est celle du noyau linéarisé sur les premières itérations utiles.

\subsection{Algorithme de calcul}

\begin{definition}[Algorithme expérimental]
Pour chaque couple \((w,c)\), on exécute :

\begin{enumerate}
    \item calcul exact de \(u^\star(w,c)\) ;
    \item itération \(u^{(r+1)}=T_\eta(u^{(r)};w,c)\) ;
    \item enregistrement des énergies \(J_{\mathrm{PCH}}(u^{(r)};w,c)\) ;
    \item mesure des ratios \(\rho_r\) uniquement lorsque \(\|u^{(r)}-u^\star\|\ge\tau_{\mathrm{fp}}\) ;
    \item stockage de l'erreur finale \(\|u^{(\mathrm{final})}-u^\star\|\) et du résidu fixe
    \[
    \|T_\eta(u^\star;w,c)-u^\star\|.
    \]
\end{enumerate}
\end{definition}

\subsection{Résultats expérimentaux}

Les quatre séries de tests ci-dessous ont été exécutées avec des dimensions et nombres de contextes différents, tout en gardant les mêmes paramètres fonctionnels
\[
\lambda_{\mathrm{ctx}}=1.7,
\qquad
\lambda_{\mathrm{ch}}=0.4,
\qquad
\lambda_{\mathrm{reg}}=0.8.
\]
Le facteur théorique correspondant est donc calculé exactement par
\[
\rho(\eta)
=
\max\left\{
\frac{1}{1+\eta(1+2\lambda_{\mathrm{tot}})},
\frac{1}{1+\eta(1+2(\lambda_{\mathrm{ctx}}+\lambda_{\mathrm{tot}}))}
\right\},
\qquad
\lambda_{\mathrm{tot}}=\lambda_{\mathrm{ch}}+\lambda_{\mathrm{reg}}=1.2.
\]

Les résultats principaux sont résumés dans le tableau suivant.

\begin{center}
\begin{tabular}{|c|c|c|c|c|c|c|}
\hline
Run & \(\eta\) & \(\rho(\eta)\) & \(\|u^{(\mathrm{final})}-u^\star\|\) moyen & \(\|T_\eta(u^\star)-u^\star\|\) moyen & ratio moyen & violations énergie \\
\hline
1 & 1.0 & 0.2273 & \(1.39\times 10^{-16}\) & \(1.12\times 10^{-16}\) & 0.3443 & 0 \\
2 & 0.5 & 0.3704 & \(1.73\times 10^{-13}\) & \(1.10\times 10^{-16}\) & 0.3583 & 0 \\
3 & 2.0 & 0.1282 & \(1.25\times 10^{-16}\) & \(1.16\times 10^{-16}\) & 0.4649 & 0 \\
4 & 1.0 & 0.2273 & \(2.31\times 10^{-16}\) & \(1.93\times 10^{-16}\) & 0.8385 & 0 \\
\hline
\end{tabular}
\end{center}

Deux faits se dégagent nettement :

\begin{itemize}
    \item le résidu au point fixe est systématiquement de l'ordre de \(10^{-16}\), ce qui valide l'exactitude algébrique du point fixe et la cohérence numérique du map proximal ;
    \item l'énergie est monotone dans toutes les expériences, sans aucune violation, ce qui confirme le caractère descendant du schéma proximal.
\end{itemize}

En revanche, le ratio de contraction moyen calculé sur toute la trajectoire n'est pas un estimateur robuste de la contraction théorique. Lorsqu'on approche du bruit machine, le quotient
\[
\frac{\|u^{(r+1)}-u^\star\|}{\|u^{(r)}-u^\star\|}
\]
devient numériquement instable. La comparaison rigoureuse doit donc se faire sur les premières itérations utiles, ou directement sur le noyau linéarisé
\[
D T_\eta(z;w,c)
=
\operatorname{diag}\!\left(
\frac{1}{1+\eta(1+2\lambda_{\mathrm{tot}})},
\frac{1}{1+\eta(1+2(\lambda_{\mathrm{ctx}}+\lambda_{\mathrm{tot}}))}
\right),
\]
dont la norme opérateur vaut exactement \(\rho(\eta)\).

Ainsi, les résultats expérimentaux montrent :
\[
u^{(r)} \to u^\star,
\qquad
J_{\mathrm{PCH}}(u^{(r)};w,c)\ \text{décroît},
\qquad
T_\eta(u^\star;w,c)=u^\star,
\]
mais ils ne constituent pas, en l'état, une validation empirique forte d'une contraction uniforme globale mesurée par un ratio moyen sur toute la trajectoire.

\subsection{Interprétation théorique}

Ces résultats valident rigoureusement le noyau proximal PCH suivant :

\begin{itemize}
    \item existence d'un unique minimiseur ;
    \item forme fermée de l'opérateur proximal ;
    \item point fixe exact ;
    \item décroissance monotone de l'énergie ;
    \item stabilité numérique au niveau machine.
\end{itemize}

Ils valident en revanche seulement de manière partielle l'assertion de contraction si celle-ci est lue via un ratio moyen global calculé sur la totalité des itérations.
La lecture correcte est la suivante : la contraction théorique est celle du noyau linéarisé, et elle doit être évaluée avant saturation numérique.

\subsection{Conséquences pour la théorie}

\paragraph{Conséquences conceptuelles.}
Cet appendice apporte une validation expérimentale nette du bloc proximal de la théorie.
Il montre que la purification contextuelle PCH peut être réalisée par un opérateur déterministe, strictement convexe, exact et stable.
Il confirme donc la cohérence mathématique du noyau de base de PCH.

\paragraph{Conséquences pratiques.}
Le protocole fournit un banc d'essai propre pour :
\begin{itemize}
    \item tester des sous-espaces contextuels réels construits à partir d'embeddings ;
    \item comparer les contractions théoriques et observées sur les premières itérations utiles ;
    \item intégrer ensuite le bloc proximal dans des architectures plus riches de type PCH-ESM ;
    \item servir de module de référence pour des expériences futures sur données réelles.
\end{itemize}

\paragraph{Portée réelle.}
Cette expérience ne valide pas encore la totalité de PCH-ESM.
Elle valide de façon solide le noyau proximal PCH, c'est-à-dire la brique élémentaire de purification contextuelle.
La généralisation à des représentations sémantiques réelles, à la décomposition spectrale complète, au collage pondéré multi-modulaire et à la dynamique Multi-PH reste un travail expérimental distinct.

\subsection{Conclusion de l'appendice}

En résumé, l'expérience confirme trois faits essentiels :
\[
\text{(i) minimiseur exact},
\qquad
\text{(ii) convergence exacte},
\qquad
\text{(iii) stabilité énergétique}.
\]
La mesure de contraction doit être interprétée avec prudence et calculée selon un protocole filtré sur les premières itérations utiles.
Sous cette lecture rigoureuse, l'appendice constitue une validation propre, fidèle et mathématiquement cohérente du noyau proximal PCH.



%%%%%%%%%%%%%%%%%%%%%%%%%%%%%%%%%%%%%%%%%%%%%%%%%%%%%%%%%%%%%%%%%%%%%%%%%%%%%%%%%%%%%%%%%%%%%%%%%%%%%%%%%%%%%%%%%%%%%%%%%%%%%%%%%%%%%%%%%%%%%%%%%%%%%%%%%%%%%%%%%%%%%%%%%%%%%%%%%%%%%%%%%%%%%%%%%%%%%%%%%%%%%%%%%%%%%%%%%%%%%%%%%%%%%%%%%%%%%%%%%%%%%%%%%%%%%%%%%%%%%%%%%%%%%%%%%%%%%%%%%%


\appendix

\section{Validation expérimentale du noyau PCH sur embeddings réels}

% =========================================================
\subsection{Positionnement de cette annexe}
% =========================================================

Cette annexe documente les expériences réalisées sur le noyau PCH dans un cadre contrôlé utilisant des représentations textuelles réelles.

L’objectif n’est pas de valider expérimentalement l’intégralité du formalisme PCH--ESM, mais de tester, dans un cadre concret, les propriétés conditionnelles établies pour le bloc PCH pris isolément : existence du minimiseur, unicité numérique, validité de l’itération proximale, contraction locale, convergence vers le point fixe, décroissance monotone de l’énergie, stabilité numérique, réduction de distance aux ensembles admissibles, et effet géométrique de la purification contextuelle sur des embeddings réels.

Le message empirique central est le suivant : le noyau PCH est stable et cohérent, mais la qualité finale du système dépend fortement du mécanisme de sélection du contexte. Autrement dit, les expériences soutiennent le bloc PCH conditionnel à un contexte donné, sans pour autant valider à elles seules le système PCH--ESM complet.

Il est important de souligner dès maintenant un point structurel majeur issu des expériences : le comportement observé est compatible avec une dynamique essentiellement quadratique locale de l’opérateur proximal, ce qui place le noyau PCH dans une classe de systèmes fortement structurés plutôt que dans une dynamique heuristique.

% =========================================================
\subsection{Question scientifique}
% =========================================================

On considère un embedding latent
\[
w \in \mathbb R^d.
\]

Étant donné un contexte \(c\), on associe un ensemble admissible contextuel \(\mathcal M_c\).

Le formalisme PCH affirme qu’il existe un unique embedding purifié
\[
u^\star = \mathcal P^{\mathrm{PCH}}(w;c),
\]
obtenu comme minimiseur d’une énergie contextuelle fortement convexe.

Les expériences cherchent à vérifier empiriquement les points suivants :
\begin{enumerate}[leftmargin=2em]
    \item l’existence numérique du minimiseur ;
    \item son unicité effective ;
    \item la validité de l’itération proximale ;
    \item la contraction annoncée par la théorie ;
    \item la convergence vers le point fixe ;
    \item la décroissance monotone de l’énergie ;
    \item la stabilité numérique du calcul ;
    \item la réduction de distance au contexte admissible ;
    \item l’impact éventuel sur une tâche aval linéaire.
\end{enumerate}

Une question secondaire, mais devenue centrale empiriquement, est de déterminer si l’opérateur proximal se comporte localement comme une application affine issue d’un problème quadratique fortement convexe.

% =========================================================
\subsection{Cadre expérimental}
% =========================================================

Les campagnes expérimentales ont été réalisées sur le corpus \texttt{20 Newsgroups} selon deux protocoles principaux :

\begin{enumerate}[leftmargin=2em]
    \item un protocole réduit à quatre contextes ;
    \item une extension à l’ensemble des vingt catégories du corpus.
\end{enumerate}

Les résultats à quatre contextes servent principalement à l’analyse du noyau PCH dans un cadre simplifié, tandis que les expériences à vingt contextes permettent d’évaluer la robustesse du mécanisme lorsque la difficulté de routage augmente.

Chaque texte est encodé par un encodeur figé \(E_{\theta_0}\) de type SentenceTransformer, puis projeté dans un espace latent de dimension finie par PCA. Aucune phase d’entraînement de l’encodeur n’est effectuée.

Les dimensions latentes explorées au cours des campagnes incluent notamment
\[
d \in \{16,32,40,45,48,50,64,80,128\}.
\]

Les tailles de calibration explorées incluent notamment
\[
n_{\mathrm{calib}} \in \{16,32,45,48,50,64,80,128\}.
\]

Les campagnes ont également été répétées sur plusieurs graines aléatoires, et les résultats les plus importants se sont montrés stables d’une graine à l’autre.

% =========================================================
\subsection{Construction des contextes admissibles}
% =========================================================

Pour chaque classe \(c\), on construit un ensemble de calibration \(D_c\). On définit ensuite le contexte admissible par
\[
\mathcal M_c = \operatorname{conv}\{E_{\theta_0}(x) : x \in D_c\}.
\]

Dans l’implémentation expérimentale, \(\mathcal M_c\) est donc une enveloppe convexe réelle dans l’espace latent. Cette construction correspond exactement à l’objet théorique utilisé dans le noyau PCH.

Les expériences confirment empiriquement que les projections sur ces ensembles jouent un rôle structurel majeur dans la stabilité globale du système.

% =========================================================
\subsection{Fonctionnelle PCH utilisée}
% =========================================================

L’énergie étudiée est
\[
J_{\mathrm{PCH}}(u;w,c)
=
\frac12\|u-w\|^2 + \Phi_{\mathrm{PCH},c}(u),
\]
avec
\[
\Phi_{\mathrm{PCH},c}(u)
=
\lambda_{\mathrm{ctx}}\operatorname{dist}(u,\mathcal M_c)^2
+
\lambda_{\mathrm{ch}}\Omega_{\mathrm{ch}}(u)
+
\lambda_{\mathrm{reg}}\|u\|^2.
\]

Dans l’implémentation expérimentale, \(\Omega_{\mathrm{ch}}(u)\) est instanciée explicitement et les sous-problèmes PCH sont résolus exactement par programmation quadratique convexe.

Aucune descente de gradient n’est utilisée, et aucune approximation stochastique n’intervient dans le calcul du prox.

% =========================================================
\subsection{Opérateur de purification}
% =========================================================

Le point purifié est défini par
\[
u^\star = \mathcal P^{\mathrm{PCH}}(w;c)
=
\arg\min_u J_{\mathrm{PCH}}(u;w,c).
\]

La théorie prévoit que ce point est bien défini et unique dès que \(\lambda_{\mathrm{reg}}>0\), ce qui est vérifié dans toutes les expériences.

Une observation importante est que ce minimiseur est atteint de manière extrêmement stable numériquement, sans instabilité de solveur ni dépendance aux conditions initiales.

% =========================================================
\subsection{Itération proximale}
% =========================================================

Le prox étudié est
\[
\mathcal T_\eta^{\mathrm{PCH}}(z;w,c)
=
\arg\min_u
\left(
J_{\mathrm{PCH}}(u;w,c)
+
\frac1{2\eta}\|u-z\|^2
\right).
\]

Les trajectoires observées sont définies par
\[
u^{(0)} = w,
\qquad
u^{(r+1)} = \mathcal T_\eta^{\mathrm{PCH}}(u^{(r)};w,c).
\]

Les valeurs testées incluent notamment
\[
\eta \in \{1,2\}.
\]

% =========================================================
\subsection{Mesures de contraction et correction méthodologique}
% =========================================================

Les premières analyses utilisaient directement les ratios
\[
\rho_r
=
\frac{\|u^{(r+1)}-u^\star\|}{\|u^{(r)}-u^\star\|}.
\]

Cette mesure devient non fiable lorsque la trajectoire entre dans le régime dominé par le bruit machine.

En pratique, dès que
\[
\|u^{(r)}-u^\star\| \lesssim 10^{-15},
\]
les fluctuations numériques dominent complètement l’estimation.

La conclusion méthodologique est donc la suivante :

\begin{itemize}[leftmargin=2em]
\item les ratios ne sont valides que dans le régime pré-asymptotique ;
\item au-delà, ils ne mesurent plus la dynamique mais l’erreur flottante ;
\item seule la phase avant domination du bruit flottant contient une information dynamique exploitable.
\end{itemize}

Cette nuance est importante : la découverte empirique n’est pas seulement que les premières itérations comptent, mais que les itérations restent informatives uniquement tant que la distance à \(u^\star\) n’a pas atteint le plancher numérique.

% =========================================================
\subsection{Linéarisation locale}
% =========================================================

L’objet pertinent pour analyser la contraction est la linéarisation locale :
\[
D\mathcal T_\eta^{\mathrm{PCH}}(u^\star;w,c).
\]

Sa norme spectrale
\[
\rho_{\mathrm{lin}}
=
\left\|D\mathcal T_\eta^{\mathrm{PCH}}(u^\star;w,c)\right\|_2
\]
a été estimée par différences finies.

Une observation majeure est la quasi-constance empirique de cette quantité sur l’ensemble des expériences, suggérant fortement une structure locale proche d’une application affine dérivée d’un problème quadratique fortement convexe.

% =========================================================
\subsection{Nature quasi-quadratique observée (résultat majeur)}
% =========================================================

Les expériences ont révélé un résultat structurel important :

\[
\rho_{\mathrm{lin}} \simeq 0.2272727
\quad \text{(très stable)}
\]

pour \(\eta=1\), et une réduction cohérente vers
\[
\rho_{\mathrm{lin}} \approx 0.128
\quad \text{pour } \eta=2.
\]

Cette stabilité forte indique que :

\begin{itemize}[leftmargin=2em]
\item l’opérateur proximal se comporte localement comme une application affine ;
\item cette application affine est compatible avec une fonctionnelle quadratique fortement convexe ;
\item la dynamique observée est donc structurellement proche d’un système linéaire contractif ;
\item le comportement n’est pas seulement contractif, il est compatible avec le prox exact d’un problème régularisé presque purement quadratique sur les domaines testés.
\end{itemize}

Autrement dit, les expériences ne montrent pas seulement une contraction, mais un comportement typique de système quadratique régularisé.

La remarquable stabilité numérique de la constante observée suggère en outre qu’elle pourrait correspondre à une expression analytique simple dépendant des paramètres du problème proximal. L’identification exacte de cette constante est laissée à un travail théorique ultérieur.

% =========================================================
\subsection{Mécanismes de routage étudiés}
% =========================================================

Plusieurs mécanismes de sélection de contexte ont été explorés.

\paragraph{Routage oracle.}
Le vrai contexte est fourni au noyau PCH. Cette configuration sert de borne supérieure conditionnelle et permet de mesurer la qualité intrinsèque du noyau indépendamment du routage.

\paragraph{Routage par énergie PCH.}
\[
c^\star_{\mathrm{energy}}(w)
=
\arg\min_{c\in\mathcal C}
J_{\mathrm{PCH}}(\mathcal P^{\mathrm{PCH}}(w;c);w,c).
\]

Cette approche est théoriquement naturelle mais empiriquement instable et bruitée dans plusieurs runs.

\paragraph{Routage par distance de projection.}
\[
c^\star_{\mathrm{proj}}(w)
=
\arg\min_{c\in\mathcal C}\operatorname{dist}(w,\mathcal M_c)^2.
\]

Cette approche est plus robuste et cohérente avec la géométrie des ensembles admissibles. Dans les expériences, elle a systématiquement mieux résisté au bruit que le routage par énergie.

L’analyse qualitative des matrices de confusion montre que les erreurs de routage ne sont pas arbitraires. Les confusions observées se produisent majoritairement entre contextes sémantiquement proches (religion, politique, matériel informatique, sports), ce qui suggère que le routage capture une structure thématique réelle malgré la difficulté accrue du problème à vingt classes.

% =========================================================
\subsection{Chemins expérimentaux explorés}
% =========================================================

Les campagnes ont exploré les configurations suivantes :
\begin{itemize}[leftmargin=2em]
    \item dimensions latentes \(d \in \{16,32,40,45,48,50,64,80,128\}\) ;
    \item tailles de calibration \(n_{\mathrm{calib}} \in \{16,32,45,48,50,64,80,128\}\) ;
    \item paramètres proximaux \(\eta \in \{1,2\}\) ;
    \item budgets d’itération \(\texttt{max\_iters} \in \{1,5,20,50,100\}\) ;
    \item différents routages (oracle, énergie, projection) ;
    \item \(\lambda_{\mathrm{ctx}} \in \{0.2, 1.7\}\).
\end{itemize}

\begin{table}[ht]
\centering
\small
\begin{tabular}{p{3.2cm}p{3.4cm}p{4.4cm}p{4.4cm}}
\toprule
Chemin & Paramètres principaux & Observation empirique & Lecture théorique \\
\midrule
Routage par énergie &
\(n_{\mathrm{calib}}=16\) à \(64\), \(d=32\) à \(64\) &
sélection de contexte souvent bruitée ; gains probe faibles ou négatifs dans plusieurs runs ; accuracies de sélection fréquemment modestes &
le noyau PCH reste stable, mais le routage n’est pas robuste \\
Routage par distance &
\(n_{\mathrm{calib}}=48\) à \(128\), \(d=48,64,128\) &
accuracies de sélection souvent autour de \(0.80\) à \(0.86\) en régime à \(4\) classes ; gains probe positifs dans plusieurs runs &
meilleure compatibilité avec la géométrie des \(\mathcal M_c\) \\
Faible dimension latente &
\(d=16\) ou \(32\) &
gains faibles ou mixtes ; distance résiduelle plus importante &
capacité géométrique insuffisante dans plusieurs régimes \\
Dimension intermédiaire ou haute &
\(d=48,64,80,128\) &
meilleure réduction de distance et meilleurs gains observés &
représentation plus expressive et séparabilité contextuelle accrue \\
Calibration plus large &
\(n_{\mathrm{calib}}=80,128\) &
sélection de contexte plus stable et meilleure réduction de distance &
les hulls \(\mathcal M_c\) deviennent plus représentatifs \\
\(\eta=2\) &
\(\eta=2\) au lieu de \(1\) &
ratio utile autour de \(0.125\), cohérent avec la linéarisation locale &
contraction plus rapide, comme attendu \\
Itérations supplémentaires &
\(\texttt{max\_iters}=1,5,20,50,100\) &
quasi pas d’effet une fois le point fixe atteint ; \(5\) itérations suffisent presque toujours &
l’information utile est concentrée dans les premiers pas \\
Extension à \(20\) classes &
\(d=48\), \(n_{\mathrm{calib}}=48\), \(\eta=1\), \(\lambda_{\mathrm{ctx}}=1.7\) &
sélection de contexte plus difficile (\(\approx 0.59\) à \(0.62\)), mais réduction de distance toujours forte et gain probe encore positif &
la difficulté du routage croît avec le nombre de contextes, mais le noyau PCH reste valide \\
\bottomrule
\end{tabular}
\caption{Chemins expérimentaux explorés pour la preuve de concept PCH.}
\label{tab:pch_paths}
\end{table}

% =========================================================
\subsection{Résultats principaux}
% =========================================================

Les expériences réalisées au cours des campagnes montrent de façon récurrente les faits suivants.

\begin{enumerate}[leftmargin=2em]
    \item \textbf{Aucune violation de décroissance d’énergie.}
    Pour tous les runs rapportés, le compteur de violations est resté nul :
    \[
    N_{\mathrm{viol}} = 0.
    \]

    \item \textbf{Résidu de point fixe à la précision machine.}
    Les résidus observés sont de l’ordre de \(10^{-17}\), ce qui est compatible avec une résolution exacte des QP convexes à l’intérieur de la précision flottante.

    \item \textbf{Contraction locale cohérente avec la linéarisation.}
    Pour \(\eta=1\), la norme linéarisée est restée très proche de \(0.2272727\), tandis que les ratios utiles mesurés sur les premières itérations étaient légèrement inférieurs.
    Pour \(\eta=2\), la norme linéarisée et les ratios utiles ont chuté vers \(\approx 0.128\), en accord avec la lecture contractive du prox.

    \item \textbf{Le routage est le goulot d’étranglement.}
    Le noyau PCH lui-même est numériquement stable, mais le gain final dépend fortement de la qualité du choix du contexte.

    \item \textbf{La distance de projection est plus robuste que l’énergie.}
    Le routage par distance de projection a fourni de meilleures accuracies de sélection de contexte que le routage par énergie dans les configurations comparées, et cela s’est traduit dans plusieurs cas par un meilleur probe après PCH.

    \item \textbf{La calibration et la dimension latente comptent.}
    Les améliorations les plus nettes apparaissent généralement quand la calibration est plus large et quand la dimension latente est suffisamment expressive.

    \item \textbf{L’augmentation du nombre d’itérations au-delà de la convergence n’apporte presque rien.}
    Les trajectoires ont déjà convergé dans les premières itérations utiles.

    \item \textbf{Le noyau PCH ne crée pas à lui seul la séparabilité.}
    Dans plusieurs runs, le probe s’améliore légèrement ; dans d’autres, il reste stable ou se dégrade légèrement. Le bloc PCH agit donc surtout comme régularisation géométrique contextuelle, pas comme générateur de structure discriminante nouvelle.

    \item \textbf{La stabilité train/test est très forte.}
    Sur les runs les plus représentatifs, les statistiques train et test sont presque identiques, ce qui suggère que le comportement observé n’est pas un artefact de partition mais une propriété structurelle du bloc PCH dans le régime testé.

    \item \textbf{L’extension à \(20\) contextes révèle une difficulté accrue du routage.}
    Quand le nombre de contextes augmente, la sélection correcte baisse nettement, alors que la contraction locale, la stabilité numérique et la réduction de distance restent intactes.
\end{enumerate}

% =========================================================
\subsection{Rapidité de convergence observée}
% =========================================================

Un résultat empirique central est la rapidité extrême de convergence.

Dans la majorité des configurations :

\begin{itemize}[leftmargin=2em]
\item \(1\) à \(5\) itérations suffisent pour atteindre le régime stationnaire ;
\item \(20\), \(50\) et \(100\) itérations n’apportent quasiment aucune amélioration supplémentaire ;
\item la dynamique utile est entièrement concentrée dans les premiers pas avant domination du bruit flottant.
\end{itemize}

Cela confirme que PCH agit comme un opérateur de projection fortement contractif plutôt que comme une dynamique lente.

% =========================================================
\subsection{Validation quantitative consolidée}
% =========================================================

Les campagnes les plus larges ont fourni des valeurs particulièrement stables.

Dans les runs à \(4\) contextes avec \(d\in\{32,48,64,80,128\}\), \(n_{\mathrm{calib}}\in\{48,64,80\}\), \(\eta=1\) et \(\lambda_{\mathrm{ctx}}\in\{0.2,1.7\}\), on observe typiquement :

\begin{itemize}[leftmargin=2em]
    \item une norme linéarisée moyenne très proche de
    \[
    \rho_{\mathrm{lin}} \approx 0.22727 ;
    \]
    \item un ratio utile moyen autour de
    \[
    \rho_{\mathrm{useful}} \approx 0.213 \text{ à } 0.224 ;
    \]
    \item des résidus moyens de point fixe de l’ordre de
    \[
    10^{-17} ;
    \]
    \item aucune violation de décroissance d’énergie ;
    \item une réduction de distance typiquement comprise entre
    \[
    0.06 \text{ et } 0.15
    \]
    selon la calibration et la dimension ;
    \item des gains probe allant de légèrement négatifs à franchement positifs selon le routage.
\end{itemize}

Dans les campagnes les plus stables à \(4\) classes, les chiffres les plus représentatifs sont les suivants :

\[
\rho_{\mathrm{useful}}^{\mathrm{train}} \approx 0.213375,
\qquad
\rho_{\mathrm{useful}}^{\mathrm{test}} \approx 0.213353,
\]
\[
\rho_{\mathrm{lin}}^{\mathrm{train}} \approx 0.227274,
\qquad
\rho_{\mathrm{lin}}^{\mathrm{test}} \approx 0.227273,
\]
\[
\text{résidu moyen de point fixe} \approx 7.8\times 10^{-17}.
\]

L’accord train/test est donc extrêmement fort : l’écart entre les ratios utiles moyens train et test est de l’ordre de
\[
|0.213375 - 0.213353| \approx 2.2\times 10^{-5}.
\]

La réduction de distance est également très marquée dans ces campagnes. Pour un régime représentatif à \(4\) classes, on observe par exemple :
\[
0.13611 \rightarrow 0.00672 \quad \text{(train)},
\qquad
0.13756 \rightarrow 0.00678 \quad \text{(test)}.
\]
Cela correspond à une réduction relative d’environ
\[
95.1\%.
\]

Pour l’un des régimes les plus favorables en \(4\) classes, avec routage par distance de projection, \(n_{\mathrm{calib}}=64\), \(d=32\), \(\eta=1\), on obtient typiquement :
\[
\text{context selection accuracy}\approx 0.84 \text{ à } 0.85,
\qquad
\text{probe gain}\approx +1.6\%.
\]

Dans un autre régime favorable, avec \(d=48\), \(n_{\mathrm{calib}}=48\), \(\eta=1\), \(\lambda_{\mathrm{ctx}}=1.7\), on observe typiquement :
\[
\text{probe gain}\approx +4\% \text{ à } +5\%,
\qquad
\Delta d \approx 0.10.
\]

Pour l’extension à \(20\) classes, avec \(d=48\), \(n_{\mathrm{calib}}=48\), \(\eta=1\), \(\lambda_{\mathrm{ctx}}=1.7\), \(\texttt{max\_iters}=5\), on observe :
\[
\text{context selection accuracy}\approx 0.59 \text{ à } 0.62,
\]
\[
\text{probe accuracy}
=
53.28\%
\rightarrow
54.06\%
\quad
(+0.78\%),
\]
\[
\text{distance reduction}\approx 0.13,
\]
tout en conservant la stabilité numérique et la décroissance d’énergie.

Ces valeurs consolident la lecture suivante : le noyau PCH est robuste, mais la sélection de contexte devient plus difficile à mesure que le nombre de contextes augmente.

% =========================================================
\subsection{Résumé quantitatif du régime à vingt contextes}
% =========================================================

Le régime à vingt contextes est particulièrement utile car il met en évidence la fragilité du routage tout en montrant que le noyau PCH reste stable.

Un résumé représentatif de ce régime est le suivant :
\begin{table}[ht]
\centering
\small
\begin{tabular}{lc}
\toprule
Mesure & Valeur \\
\midrule
Context selection accuracy (train) & 0.6234 \\
Context selection accuracy (test) & 0.5859 \\
Distance avant (train) & 0.1361 \\
Distance après (train) & 0.0067 \\
Distance avant (test) & 0.1376 \\
Distance après (test) & 0.0068 \\
Probe brut & 53.28\% \\
Probe après PCH & 54.06\% \\
Gain probe & +0.78\% \\
Résidu moyen & $7.8\times10^{-17}$ \\
$\rho_{\mathrm{lin}}$ & 0.22727 \\
\bottomrule
\end{tabular}
\caption{Résumé du régime représentatif à vingt contextes.}
\label{tab:twenty_context_summary}
\end{table}

% =========================================================
\subsection{Hiérarchie empirique des facteurs}
% =========================================================

Les expériences permettent d’établir la hiérarchie suivante des facteurs influençant les performances :

\[
\text{Routage}
\gg
d
\approx
n_{\mathrm{calib}}
\gg
\eta
\gg
\text{nombre d’itérations}.
\]

Cette hiérarchie est robuste sur l’ensemble des runs. Elle est encore plus nette lorsque l’on compare les régimes à \(4\) contextes et à \(20\) contextes : l’augmentation du nombre de contextes dégrade principalement le routage, sans affecter la validité interne du noyau PCH.

% =========================================================
\subsection{Absence de création de séparabilité}
% =========================================================

Une observation importante est la suivante :

\begin{itemize}[leftmargin=2em]
\item PCH améliore parfois les performances du probe ;
\item PCH peut aussi ne produire aucun gain ;
\item PCH ne génère pas systématiquement une meilleure séparabilité ;
\item PCH peut même légèrement dégrader la performance aval lorsqu’une partie du routage est mal alignée avec la structure sémantique.
\end{itemize}

Ainsi :
\[
\text{PCH ne crée pas la structure discriminante}
\]
mais agit comme un opérateur de régularisation géométrique conditionnelle.

Il exploite la structure existante du latent sans en inventer une nouvelle. Cette conclusion est importante : le bloc PCH est un mécanisme de purification contextuelle, pas un générateur universel de séparabilité de classes.

% =========================================================
\subsection{Synthèse empirique consolidée}
% =========================================================

La synthèse de l’ensemble des runs est la suivante.

\begin{itemize}[leftmargin=2em]
    \item Le noyau PCH est \emph{stable} et \emph{numériquement cohérent}.
    \item La \emph{contraction locale} attendue est observée de manière récurrente.
    \item Les \emph{résidus de point fixe} sont systématiquement au voisinage de la précision machine.
    \item La \emph{décroissance d’énergie} n’a jamais été violée dans les runs observés.
    \item Le \emph{gain aval} sur le probe linéaire est \emph{dépendant du routage}, de la calibration et de la dimension latente.
    \item Le routage par \emph{distance de projection} est plus robuste que le routage par énergie.
    \item L’augmentation de \(n_{\mathrm{calib}}\) améliore la stabilité contextuelle.
    \item Une dimension latente trop faible peut limiter la qualité géométrique du contexte.
    \item La bonne lecture des ratios de contraction exige de se limiter aux itérations avant domination du bruit flottant.
    \item Le comportement du prox est compatible avec une dynamique locale quasi quadratique.
    \item La stabilité train/test est très forte, ce qui plaide contre une interprétation en simple artefact numérique.
    \item L’extension à \(20\) contextes confirme la validité interne du noyau, tout en rendant le routage nettement plus difficile.
    \item Les confusions de routage observées à vingt classes sont majoritairement sémantiquement proches, ce qui suggère que le mécanisme capture une structure thématique réelle plutôt qu’un bruit aléatoire.
\end{itemize}

% =========================================================
\subsection{Interprétation}
% =========================================================

Les expériences valident fortement le noyau PCH conditionnel.

Elles montrent que :

\begin{itemize}[leftmargin=2em]
\item le problème proximal est bien posé ;
\item la dynamique est contractive ;
\item la structure locale est fortement compatible avec une dynamique quadratique fortement convexe ;
\item la projection vers les ensembles admissibles est effective ;
\item la stabilité numérique est excellente ;
\item la réduction de distance aux ensembles admissibles est forte et répétable ;
\item la stabilité train/test est très élevée ;
\item l’effet du noyau est essentiellement géométrique et non discriminant par lui-même ;
\item la réduction de distance atteint des niveaux très élevés dans les régimes les plus stables, avec une baisse typique de l’ordre de \(95\%\) dans certaines campagnes représentatives.
\end{itemize}

Cependant, elles montrent également que :

\begin{itemize}[leftmargin=2em]
\item le routage de contexte est un problème séparé ;
\item il n’est pas résolu par la théorie PCH elle-même ;
\item il est le principal facteur limitant du système global ;
\item sa difficulté croît avec le nombre de contextes.
\end{itemize}

% =========================================================
\subsection{Statut exact de validation (résultat crucial)}
% =========================================================

Le statut correct de validation est le suivant :

\[
\boxed{\text{Les expériences valident PCH conditionnel à } c}
\]

mais ne valident pas :

\[
\boxed{\text{le routage}}
\quad \text{ni} \quad
\boxed{\text{PCH-ESM complet}}.
\]

Cette distinction est essentielle pour toute interprétation scientifique correcte du système.

Autrement dit, les expériences confirment le noyau proximal conditionnel, mais elles ne fournissent pas encore une théorie complète du choix automatique de contexte.

% =========================================================
\subsection{Conséquences pour la théorie}
% =========================================================

Les conséquences sont les suivantes :

\begin{enumerate}[leftmargin=2em]
\item PCH est un opérateur proximal cohérent et stable ;
\item la structure locale observée est fortement compatible avec une dynamique quadratique fortement convexe ;
\item le routage doit être formalisé séparément ;
\item les extensions du système doivent être analysées indépendamment ;
\item la validation conditionnelle ne se propage pas automatiquement au système global ;
\item l’augmentation du nombre de contextes fait apparaître un problème de routage plus difficile, sans remettre en cause la validité interne du noyau ;
\item la mesure de contraction doit être lue dans le régime avant bruit machine ;
\item la stabilité train/test très forte suggère que les phénomènes observés sont structurels et non accidentels ;
\item la réduction de distance aux ensembles admissibles constitue un effet géométrique robuste du noyau ;
\item le noyau PCH ne doit pas être présenté comme un générateur universel de séparabilité ;
\item à vingt contextes, les erreurs de routage s’organisent surtout entre catégories thématiquement proches, ce qui donne une interprétation géométrique et sémantique aux confusions observées.
\end{enumerate}

% =========================================================
\subsection{Conséquences conceptuelles, opérationnelles et pratiques}
% =========================================================

Sur le plan conceptuel, les résultats soutiennent l’idée qu’une représentation utile n’est pas seulement prédictive, mais aussi stable sous purification contextuelle. Ils soutiennent aussi l’idée qu’un espace sémantique monolithique peut être remplacé par une structure conditionnée par le contexte.

Sur le plan opérationnel, le meilleur protocole observé est :
\[
\text{encodeur figé}
\;\to\;
\text{PCA}
\;\to\;
\text{routage géométrique}
\;\to\;
\text{PCH exact}
\;\to\;
\text{probe linéaire}.
\]

Dans ce cadre, les gains les plus crédibles apparaissent avec des dimensions latentes modérément élevées à élevées et avec une calibration plus large.

Sur le plan pratique, l’expérience indique qu’il faut d’abord stabiliser le routage, puis augmenter la capacité géométrique du latent, puis seulement comparer les gains de classification. Autrement dit, le PCH exact est déjà exploitable, mais son efficacité dépend encore de la qualité du mécanisme de sélection de contexte.

% =========================================================
\subsection{Comparaison avec l’état de l’art}
% =========================================================

Cette expérience n’est pas une comparaison SOTA au sens strict. Elle utilise un encodeur figé, un sous-ensemble de \texttt{20 Newsgroups}, et un probe linéaire. La bonne comparaison est donc interne : embeddings bruts contre embeddings purifiés par PCH, sous le même protocole.

Sous ce protocole, PCH produit parfois des gains nets et parfois des gains nuls ou légèrement négatifs, selon le routage, la dimension latente et la calibration. Il faut donc présenter ces résultats comme une preuve de concept du mécanisme PCH, et non comme une domination générale de l’état de l’art.

L’intérêt scientifique de l’expérience n’est pas de battre un modèle spécialisé, mais de montrer que le bloc PCH possède une dynamique interne propre, stable, analytiquement lisible et expérimentalement reproductible.

% =========================================================
\subsection{Conclusion générale}
% =========================================================

L’ensemble des expériences conduit à une conclusion robuste :

\begin{itemize}[leftmargin=2em]
\item le noyau PCH est stable et bien défini ;
\item le minimiseur existe et est unique numériquement ;
\item le prox converge rapidement ;
\item l’énergie décroît systématiquement ;
\item la contraction locale est confirmée ;
\item la structure observée est compatible avec un système quasi-quadratique ;
\item la réduction de distance aux ensembles admissibles est forte et répétable ;
\item la stabilité train/test est très élevée ;
\item la principale limite du système est le routage ;
\item l’augmentation du nombre de contextes rend le routage plus difficile, sans invalider le noyau ;
\item le noyau PCH ne crée pas à lui seul la séparabilité ;
\item l’effet du noyau est d’abord géométrique, pas discriminant au sens supervisé.
\end{itemize}

Ainsi, les expériences apportent une validation solide du noyau PCH conditionnel, tout en identifiant clairement le routage comme le problème ouvert fondamental pour l’extension vers un système PCH--ESM complet.
\end{document}